% NAL 1644, batch 7 — Gallica views 121–140.
% Pass: Opus 5.5 (claude-opus-5-5), 2026-09-23 — first pass, unchecked against
% the views by a human.
% Read from the local mirror only (archives/nal-1644/f0121.jpg … f0140.jpg):
% a strip of thumbnails of the twenty views, a 1100-px view of each, then four
% to six full-width bands at 2400 px (normalised), and tight crops at 1.5–3×
% for marginal notes, struck words, corrected signs and numbers. Each view was
% drafted as soon as it was read.
%
% Dating. The catalogue gives no dating for this volume: \dating{} is left
% empty, and this pass infers no date. No date is written on these leaves.
%
% What the views are. Greyscale scans, portrait, about 3750 × 5620, one page
% per view. As in batch 6, the leaves are of unequal height and the camera
% often caught two at once: the top lines of a taller leaf above a shorter one
% (views 122, 124, 130, 132, 138, 140), or the foot of the leaf beneath below
% the edge of the page (views 126, 127, 128, 129, 134, 135, 136, 137), or a
% strip of a neighbouring leaf at one side. Such overlaps are transcribed once,
% on the view where their page is seen whole, and a \note{} at the head of each
% view says what else shows. Every one of the twenty views carries text; none
% is skipped.
%
% What the twenty views hold. One continuous run of Latin algebra in the
% notation of Viète's school, in the same hand as batch 6, in chapters
% numbered by the writer, each proposition followed by numerical examples in
% cossic signs (N, Q, C, QQ, QC, L, LC, LVC).
%
%   v. 121       End of a problem begun before this batch (a cubic with the
%                cube affected by the side added, solved by two « latera » D
%                and E; examples 1C + 81N = 702, 1C + 6N = 2), then
%                « Problema II » (cube affected by a negated side, reduced to
%                a « plane under a solid root negated from the square »).
%   v. 122       End of Problema II, its « Consectarium » (A = LVC(Z + L(Z² −
%                B³)) + LVC(Z − L(Z² − B³))), example 1C − 81N = 756, then
%                the title and first paragraph of « De Canonica Æquationum
%                transmutatione … Cap VIII ».
%   v. 123–125   Cap. VIII continued: numerical tests for rational roots
%                (1C ± 1N), then the canonical transmutation of the
%                coefficient or of the homogeneum of comparison, three worked
%                examples (1C + 20Q = 96,000; 1C − 144N = 10368;
%                1C + 860N = 1,728).
%   v. 125–129   « Anomala æquationum aliquot cubicarum ad quadraticas aut
%                etiam simpliciores reductio. Cap IX. », propositions 1 to
%                XI, each cubic depressed to a quadratic by a known root
%                (division by A ± B), with examples; prop. IIII only in a
%                marginal note at the foot of v. 126; props. X–XI with a
%                verification calculation at the foot of v. 129.
%   v. 130–134   « Similium reductionum Continuatio Cap. XX » (corrected in
%                the margin to X), propositions 1 to VI; prop. VI is a
%                sextic (« quadrato-quadratica ») reduced through a square.
%   v. 135–136   « Singularium aliquot constitutionum ad æqualitates
%                multipliciter affectas pertinentium Collectio. Cap. XXI »
%                (margin: XI), props. 1–V: A² + BA = B², A³ + BA² + B²A = B³,
%                … read as the division of the double of B in extreme and mean
%                ratio « duplicata, triplicata … quintuplicata ».
%   v. 136–137   « Earundem collectio altera Cap. XXII » (margin: XII),
%                props. 1–V: B the first of a series of continued
%                proportionals, Z the sum of the others, A the second.
%   v. 137–139   « Earundem collectio tertia Cap. XXIII » (margin: XIII),
%                props. 1–IIII: the same with alternating differences; five
%                further propositions of the series in marginal notes of
%                v. 137 and 138.
%   v. 139–140   « Collectio quarta. Cap. XXIIII » (margin: XIIII), on
%                equations « de multiplicibus radicibus mire explicabiles »:
%                prop. 1 (two roots B, D), II (three), III (four: example
%                1, 2, 3, 4), IIII (five roots), whose statement runs off the
%                foot of v. 140.
%
%   The chapter numbers are written XX, XXI, XXII, XXIII, XXIIII in the text
%   and corrected, each time, to X, XI, XII, XIII, XIIII in the left margin,
%   in continuation of « Cap VIII » (v. 122) and « Cap IX » (v. 125). The
%   leaves name no author and carry no title for the whole; the pass
%   attributes nothing.
%
% Continuity. View 121 continues a text begun before it: it opens on the last
% equality of a solution whose problem is not on the view (the end of
% batch 6, or a leaf before it). The text then runs without a break from view
% to view, recto to verso and leaf to leaf, to view 140, whose last
% proposition (Cap. XIIII, prop. IIII) stops mid-statement at the foot of the
% page, its brace open and no second member written: the batch runs on into
% view 141.
%
% Hands. The text is one regular cursive throughout. A second layer, in a
% blacker ink and a more pressed stroke, corrects it: underlined words or
% signs in the text with the correct form written, underlined, in the left
% margin (a coefficient, a sign, a whole term); longer marginal notes that
% justify a rule (v. 122) or state further propositions of a series
% (v. 126, 134, 137, 138); and a boxed note at v. 138. Whether this layer is
% the same writer at a second sitting or another reader, the pass does not
% say. Marginal corrections are given as \marginal{} beside the line they
% answer, and a \note{} says what they replace; the text is left as written.
%
% Notation. Viète's notation stays as written: the unknown A (E for a new
% unknown), the given magnitudes B, D, G, H, K, X, Z named by their dimension
% (« B plano », « Z solido », « B plano-plano-plano »), the homogeneity kept,
% « in » for a product, « æquatur / æquabitur » for the equality, braces
% gathering the terms of a member. Abbreviated powers in the margins (Aq, Ac,
% Aqq, Bpl., Bppp, Xqq, Zs., Zss.) are kept as written. The numerical examples
% use the cossic signs N, Q, C, QQ, QC, with L for a square root, LC for a cube
% root, LVC for the cube root of a whole binomial. Terms written in columns are
% set as arrays in the page's column order. A double stroke « = » stands
% before a term where the writer's sign was retraced or left undecided; it is
% not an equals sign and is kept as it stands. A cross laid over a double
% stroke in the corrector's margin is set $\pm$. Roman numerals I, II, III…
% written as strokes over a row of continued proportionals are set with
% \overset. The long s is set s.
%
% Numbers on the leaves. One ink figure at the top right of each recto, in
% the series batch 6 reads up to 56 (v. 120): 57 (v. 122), 58 (v. 124),
% 59 (v. 126, and above the slips at v. 122 and 124), 60 (v. 128), 61
% (v. 130), 62 (v. 132), 63 (v. 134, and above the leaves at v. 130 and 132),
% 64 (v. 136), 65 (v. 138), 66. (v. 140), and 67 on the leaf whose top shows
% above v. 138 and 140. The figure shows through, mirrored, at the top left of
% several versos (v. 127, 129, 131, 139); the versos carry no number of their
% own. One figure per leaf, running on without a gap: a leaf count, not the
% writer's pagination. It is ink, not a pale pencil series, and no library
% certificate is in this batch, so no \folio{} is written. The page of
% view 121 carries no number; by the series it would be the verso of leaf 56,
% which the pass does not assert. Other numbers are the writer's: chapter and
% proposition numbers, the corrector's marginal figures, the calculations.
%
% Other signs. No stamp, no shelfmark, no date on these twenty views.

% Coordinator's reconciliation (2026-09-23), after all nine batches:
%   One ink series of leaf numbers runs top right of the rectos through the
%   whole volume, with no gap: 1–8 and « 5 bis » (v. 5–20), 9–17 (v. 22–39),
%   18–27 (v. 41–59), 28–37 (v. 61–79), 38–47 and « 46 bis » (v. 81–100),
%   « 47 bis » and 48–56 (v. 102–120), 57–66 (v. 122–140), 67–78
%   (v. 142–160), 79 (v. 162). It crosses the three pieces of the volume, and
%   the writer cites it himself (« fol. 18. », v. 43; « pag. 48 », v. 100), so
%   it is the writer's or copyist's count of leaves, not the library's. Being
%   ink, it gets no \folio{}, in any batch.
%   The pieces: « Francisci Vietæ De Recognitione Æquationum tractatus »
%   (heading on v. 5), Cap. 1–XXI, v. 5–68; « Francisci Vietæ. De æquationum
%   Emendatione tractatus secundus » (heading on v. 69), Cap. I–XIIII, ending
%   « Finis » on v. 141; then astronomy without a name, v. 142–163 (lunar
%   prostaphaereses, Tycho's lunar hypothesis restated, « Harmonici
%   Ptolemaici Constructio Caput XII », Mercury). Two doubtful notes may mark
%   the treatises as a copy: « transcripta ex exemplari » (v. 100) and
%   « deerant hæc in exemplari » (v. 114). The name on the leaves is in the
%   headings; who wrote the hands is not settled.

\documentclass[11pt,a4paper]{article}
% The preamble shared by every transcription of the mathematicians' manuscripts in Gallica.
%
% It rests on one idea: **uncertainty compounds, it does not resolve.** A
% transcription that smooths over a doubtful word has destroyed the only thing
% it was made for — a reader must be able to tell, at every point, what was on
% the page and what was supplied. The apparatus macros below are therefore the
% critical apparatus, not decoration.
%
% The same commands are recognised by `scripts/render.mjs`, which produces the
% site's reading view, and by `scripts/tei.mjs`. A macro added here must be
% added there too, or rendering fails loudly — which is the wanted behaviour.
%
% Comments here are English, like the rest of the repository. The documents
% this preamble sets are French, and that is deliberate: see the skills.

\ProvidesPackage{germain}[2026/09/15 Transcriptions of mathematicians' manuscripts in Gallica]

\usepackage{amsmath,amssymb,amsthm}
\usepackage{mathrsfs}
% Kept from the parent preamble so the renderer's diagram support stays
% available; these manuscripts draw no commutative diagrams, and nothing here uses it.
\usepackage{tikz-cd}
\usepackage[english,french]{babel}
\usepackage{fontspec}

% --- Greek ------------------------------------------------------------------
% The text face stays fontspec's default, Latin Modern, which has no Greek:
% XeTeX drops every glyph it lacks with a line in the log and nothing on the
% page, and the Greek words the manuscripts quote vanished from the PDFs. So
% the Greek and Greek Extended blocks get a XeTeX character class of their own,
% and entering or leaving a run of it switches the family to CMU Serif — the
% same Computer Modern design, with polytonic Greek — and back. Only the
% family changes: a Greek word in \textit or \textbf keeps its shape and series.
% The transitions to and from every other class carry the tokens that class
% has towards an ordinary letter, so French spacing before « ; » or inside
% « » still holds after a Greek word. No grouping, so no brace or \endgroup
% can come out unbalanced; maths is left alone, since XeTeX inserts nothing
% there. Kept identical in `exercices.sty`.
\makeatletter
\newfontfamily\germainGreekfont{cmun}[
  Extension=.otf, NFSSFamily=germaingreek,
  UprightFont=*rm, ItalicFont=*ti, SlantedFont=*sl,
  BoldFont=*bx, BoldItalicFont=*bi, BoldSlantedFont=*bl]
\newXeTeXintercharclass\germain@greek
\def\germain@greekrange#1#2{%
  \@tempcnta=#1\relax
  \loop
    \XeTeXcharclass\@tempcnta=\germain@greek
    \ifnum\@tempcnta<#2\advance\@tempcnta\@ne\repeat}
\germain@greekrange{"0370}{"03FF}
\germain@greekrange{"1F00}{"1FFF}
\def\germain@greekprev{\rmdefault}
\def\germain@greekin{%
  \edef\germain@greekprev{\f@family}\fontfamily{germaingreek}\selectfont}
\def\germain@greekout{\fontfamily{\germain@greekprev}\selectfont}
\def\germain@greekjoin{%
  \XeTeXinterchartokenstate=\@ne
  \@tempcnta=\z@
  \loop
    \ifnum\@tempcnta=\germain@greek\else
      \XeTeXinterchartoks\germain@greek\@tempcnta=\expandafter{%
        \expandafter\germain@greekout\the\XeTeXinterchartoks\z@\@tempcnta}%
      \XeTeXinterchartoks\@tempcnta\germain@greek=\expandafter{%
        \the\XeTeXinterchartoks\@tempcnta\z@\germain@greekin}%
    \fi
    \ifnum\@tempcnta<4095\advance\@tempcnta\@ne\repeat}
% babel-french sets its own classes, and saves and restores the interchar
% state, each time a language is selected: join again after it does.
\addto\extrasfrench{\germain@greekjoin}
\addto\extrasenglish{\germain@greekjoin}
\AtBeginDocument{\germain@greekjoin}
\makeatother

\usepackage{geometry}
\geometry{a4paper,margin=25mm,marginparwidth=28mm}
\usepackage{marginnote}
\usepackage{xcolor}
\usepackage[normalem]{ulem}
\setlength{\emergencystretch}{3em}
\usepackage{hyperref}
\hypersetup{colorlinks=true,linkcolor=black,urlcolor=black,pdfborder={0 0 0}}

% --- Batch metadata ---------------------------------------------------------
% Taken as they stand by the rendering script, which shows them at the head of
% the reading view. They fix what the file claims to cover. The macro names are
% the parent project's, so that the scripts stay shared; \folder here means the
% volume's slug — `fr-9115` — and \pages counts Gallica views.

\newcommand{\folder}[1]{\gdef\grothFolder{#1}}
\newcommand{\batch}[1]{\gdef\grothBatch{#1}}
\newcommand{\pages}[2]{\gdef\grothFirst{#1}\gdef\grothLast{#2}}
\newcommand{\foldertitle}[1]{\gdef\grothFoldertitle{#1}}

% The holder's own shelfmark, printed as they write it: « Français 9115 ».
\newcommand{\shelfmark}[1]{\gdef\grothShelfmark{#1}}

% The dating, copied verbatim from the holder's catalogue — never inferred
% here. « XIXe siècle » is the BnF's; a pass that dates a leaf from its content
% says so in a \note{}, as its own inference.
\newcommand{\dating}[1]{\gdef\grothDating{#1}}

% The status notice, printed in the title block and repeated as a diagonal
% watermark on every page of the PDF. The manuscripts are in the public
% domain; what this line declares is not a right but a quality — an
% unverified first pass by a machine. The skills write the line; a file
% without it gets no watermark, so leaving it out is visible in the output.
\newcommand{\watermark}[1]{\gdef\grothWatermark{#1}}

\gdef\grothFolder{?}\gdef\grothBatch{}\gdef\grothFirst{?}\gdef\grothLast{?}%
\gdef\grothFoldertitle{}\gdef\grothShelfmark{}\gdef\grothDating{}\gdef\grothWatermark{}

% --- The résumé -------------------------------------------------------------
\newenvironment{resume}
  {\small\setlength{\parskip}{0.45em}}
  {\par\medskip\hrule\medskip}

% --- Keywords ---------------------------------------------------------------
% In English, closing the modernised reading's résumé: the modern vocabulary
% under which to search for what the leaves construct. The manifest extracts
% them to tag the volume — this line is the single source of the tags.
\newcommand{\keywords}[1]{\par\smallskip\noindent
  {\footnotesize\textsc{Keywords} --- \textit{#1}}\par}

% --- The critical apparatus -------------------------------------------------

% The Gallica view number, in the margin. This is the anchor that lets a
% disagreement over a formula be settled by looking at *one* image rather than
% a volume of 750. The site uses it to turn the facsimile as the transcription
% is scrolled. A view is one scanned image: usually one side of a leaf.
\newcommand{\page}[1]{%
  \marginnote{\footnotesize\color{gray!70}\textsf{v.\,#1}}%
  \label{page:#1}\ignorespaces}

% The BnF's pencilled foliation, where the leaf carries it and a pass can read
% it: « 198r ». This is what the literature cites — « ff. 198r–208v » — and
% the one line that turns a citation into a view. Recorded, never inferred:
% a folio a pass cannot read is not written.
\newcommand{\folio}[1]{%
  \marginnote{\footnotesize\color{gray!50}\textsf{f.\,#1}}\ignorespaces}

% The modernised reading's coarser counterpart to \page{}: which views a
% section is drawn from.
\newcommand{\pagerange}[2]{%
  \marginnote{\footnotesize\color{gray!70}\textsf{v.\,#1--#2}}%
  \ignorespaces}

% Illegible: never guessed.
\newcommand{\ill}{\textcolor{red!60!black}{[\,\ldots\,]}}

% A doubtful reading: the word is given, and flagged as uncertain.
\newcommand{\uncertain}[1]{\uline{#1}}

% An editorial addition — a missing letter, an implied word.
\newcommand{\add}[1]{\textcolor{blue!50!black}{[#1]}}

% Struck out by the writer. What was crossed out is part of the document.
\newcommand{\struck}[1]{\sout{#1}}

% The transcriber's note, on the state of the page or the mathematics.
\newcommand{\note}[1]{\par\smallskip{\small\color{gray!60!black}\textit{[#1]}}\par\smallskip}

% A marginal note of hers, distinct from the body of the text.
\newcommand{\marginal}[1]{\par\smallskip\noindent\hspace{1em}%
  {\small\color{gray!60!black}#1}\par\smallskip}

% --- Title ------------------------------------------------------------------
% The title block is French, like the documents themselves.

\AddToHook{shipout/background}{%
  \ifx\grothWatermark\empty\else
    \begin{tikzpicture}[remember picture, overlay]
      \node[rotate=52, scale=3.4, align=center, text=gray!18,
            font=\sffamily\bfseries]
        at (current page.center) {\grothWatermark};
    \end{tikzpicture}%
  \fi}

\AtBeginDocument{%
  \begin{center}
    {\large\bfseries Manuscrits de mathématiciens (Gallica) ---
      \ifx\grothShelfmark\empty\grothFolder\else\grothShelfmark\fi}\\[2pt]
    {\itshape\grothFoldertitle}\\[4pt]
    {\small vues \grothFirst--\grothLast
      \ifx\grothBatch\empty\else\ (lot \grothBatch)\fi}\\[2pt]
    \ifx\grothDating\empty\else
      {\small\color{gray!60!black}Datation du catalogue : \grothDating}\\[2pt]
    \fi
    \ifx\grothWatermark\empty\else
      {\small\bfseries\color{red!45!gray}\grothWatermark}\\[2pt]
    \fi
    {\footnotesize\color{gray!60!black}Transcription établie d'après les images IIIF de Gallica.
     Source gallica.bnf.fr / Bibliothèque nationale de France.\\
     Les lectures incertaines sont soulignées, les illisibles notées [\,\ldots\,],
     les ajouts entre crochets ; \textsf{v.} numérote les vues de Gallica,
     \textsf{f.} la foliotation au crayon de la BnF.}
  \end{center}
  \vspace{1em}}

\endinput


\folder{nal-1644}
\shelfmark{NAL 1644}
\batch{7}
\pages{121}{140}
\dating{}
\watermark{Lecture automatique\\non vérifiée}
\foldertitle{Viète (François). Papiers divers. NAL 1644}

\begin{document}

\page{121}
\note{Page sans numéro en tête ; elle s'ouvre sur la fin d'une résolution dont l'énoncé n'est pas sur la vue, et continue donc un texte commencé avant la vue 121. Le bord supérieur du feuillet est abîmé ; à droite, le bord d'un onglet de montage.}

\marginal{$\dfrac{Eq. - Bp.}{E} \;/\; A$}
\note{Dans la marge de gauche, en tête : la même égalité que la ligne qui suit, en abrégé, la fraction séparée de « A » par un long trait oblique.}

\[ \left.\begin{array}{l} \overline{E\ \text{quadratum}} \\ -\,B\ \text{plano}\ \struck{E} \end{array}\right\}\ \text{æquabitur } A \]
\note{Un trait passe au-dessus de « E quadratum » et un autre sous « B plano » ; aucun dénominateur n'est écrit sous le second, et la lettre qui suit « B plano » est barrée. La note marginale donne $E$ pour dénominateur.}

Si $1C + 81N$ æquetur $702$

quoniam $378$ plus $351$ est $729$. Cubus a latere $9$ ideo $\dfrac{81 - 27}{9}$ seu $6$. est $1N$ de qua quæritur.
\note{« 81 » est écrit au-dessus de « $-27$ », le tout sur un trait, « 9 » dessous. Le calcul est juste : $351^2 + 27^3 = 142884 = 378^2$, $378 + 351 = 729 = 9^3$, et $9 - 27/9 = 6$, qui satisfait $6^3 + 81 \cdot 6 = 702$.}

\uncertain{Consectarium} \uncertain{ex} nobis antedictis consectarijs

Deniq; sunt duo latera unum idemque minus $D$ alterum idemque maius $E$. quorum differentia est $A$ de qua quæritur
\note{Le premier mot de l'alinéa se lit « Conſectarum » ou « Consectarium » ; le mot suivant, bref et pointé, « ex » ou « a ». La lettre lue $E$ a ici, comme partout sur la page, la forme d'un G ; celle lue $D$ est nette.}

\[ \text{Itaque } LVC.\ LQ \left\{\begin{array}{l} B\ \text{plano-plano-plani} \\ +\,Z\ \text{solido-solido} \end{array}\right. +\,Z\ \text{solido} \]
\[ -\,LVC.\ LQ \left\{\begin{array}{l} B\ \text{plano-plano-plani} \\ +\,Z\ \text{solido-solido} \end{array}\right. -\,Z\ \text{solido} \]
est $A$ quæsita
\note{« LVC » : latus universale cubicum ; « LQ » : latus quadratum. Deux lettres de « quæsita » sont surchargées.}

Si $1C + 6N$ æquetur $2$

$LC\,4 - LC\,2$ est $1N$ de qua quæritur.
\note{L'exemple est juste : pour $x^3 + 6x = 2$, $\sqrt[3]{4} - \sqrt[3]{2}$.}

\subsection*{Problema II.}

Cubum adfectum sub latere negato ad planum sub radice solida negatum de quadrato reducere.

Oportebit autem in æquatione proposita cubum \struck{\ill{}} \uncertain{de} triente coefficientis adfectionis cedere solidi comparationis subquadruplo quadrato.
\note{Le mot barré en fin de ligne est noirci. « cedere » est lu d'après les lettres (« cedere ») ; la condition énoncée est que le cube du tiers du coefficient de l'affection soit moindre que le quart du carré du solide de comparaison.}

\[ \text{Proponatur } \left.\begin{array}{l} A\ \text{cubus} \\ -\,B\ \text{plano ter in } A \end{array}\right\}\ \text{æquari } Z\ \text{solido bis} \]
oportet facere quod propositum est
\[ \left.\begin{array}{l} A \text{ in } E \\ -\,E\ \text{quadrato} \end{array}\right\}\ \text{æquetur } B\ \text{plano} \]
Unde $B$ planum ex huiusmodi æquationis \struck{compositione} intelligitur rectangulum sub duobus lateribus quorum minus est $E$. \struck{summa} vero maioris ac minoris
\marginal{constitutione.}
\note{« compositione » est barré, et « constitutione » écrit en regard dans la marge de gauche, qui le remplace. « summa » est barré sans remplaçant sur la ligne ; la ligne suivante commence par un « A » écrit au-dessus de « Igitur », qui achève la phrase : la somme du plus grand et du plus petit côté est $A$.}

$A$

\[ \text{Igitur } \dfrac{\begin{array}{l} B\ \text{plano} \\ +\,E\ \text{quadrato} \end{array}}{E} \Big\}\ \text{æquabitur } A \]
\[ \text{quare } \left.\begin{array}{l} \dfrac{\begin{array}{l} B\ \text{plano-plano-plano} \\ +\,E\ \text{quad. in } B\ \text{plano-plano ter} \\ +\,\struck{\ill{}}\,E\ \text{quad.-quad. in } B\ \text{plano ter} \\ +\,E\ \text{cubo-cubo} \end{array}}{E\ \text{cubo}} \\[2ex] \dfrac{\begin{array}{l} -\,B\ \text{plano-plano ter} \\ -\,B\ \text{plano in } E\ \text{quadrat. ter} \end{array}}{E} \end{array}\right\}\ \text{æquabitur } Z\ \text{solido} \]
\marginal{bis}
\note{« Z solido » est souligné, et « bis », souligné aussi, est écrit en regard au bas de la marge de gauche : $Z$ solido bis, comme dans l'énoncé. Devant « E quad.-quad. », un petit mot barré. Le développement est juste : c'est $A^3 - 3BA$ pour $A = (B + E^2)/E$. La page s'arrête sur cette égalité ; le bas du feuillet est blanc.}

\page{122}
\note{La vue montre deux feuillets superposés. En haut, sur une bande de deux lignes, le haut d’un autre feuillet, numéroté « 59 » à l'encre dans l'angle supérieur droit (« omnia applicentur ad $A + B$ » ; « Et oritur $A$ quadrato $- B$ in $A$ $\}$ fit $B$ quadratum bis ») : c'est le feuillet vu en entier à la vue 126, où il est transcrit. Posé dessus, un feuillet plus petit, numéroté « 57 » à l'encre en haut à droite, dont on transcrit ici la page. Elle enchaîne sur la dernière égalité de la vue 121, qu'elle multiplie par $E$ cube.}

Et omnibus per $E$ cubum ductis et ex arte concinnatis
\[ \left.\begin{array}{l} Z\ \text{solidum bis in } E\ \text{cubum} \\ -\,E\ \text{cubi quadrato} \end{array}\right\}\ \text{æquabitur } B\ \text{plani cubo} \]
quæ æquatio plani sub radice solida negati de quadrato facta, itaque est reductio quæ imperabatur

apparet autem ex æquatione illius ad quam reductio facta est proprietate $Z$ solidi quadrat. præstare debere $B$ plani cubo. quo pertinet apposita huic problemati

\subsection*{Consectarium.}

\[ \text{Itaque si } \left.\begin{array}{l} A\ \text{cubus} \\ -\,B\ \text{plano ter in } A \end{array}\right\}\ \text{æquetur } Z\ \text{solido bis} \]
\[ LVC.\ Z\ \text{solidi} + L \left\{\begin{array}{l} Z\ \text{solido-solidi} \\ -\,B\ \text{plano-plano-plano} \end{array}\right. \]
\[ +\,LVC\ Z\ \text{solidi} - L \left\{\begin{array}{l} Z\ \text{solido-solidi} \\ -\,B\ \text{plano-plano-plano} \end{array}\right. \]
est $A$ de qua quæritur
\note{Deux taches d'encre sur les racines, sans effacer de lettre.}

\marginal{quia enim $Z$ solid. bis tanquam summa fit ex \ill{} binis \ill{} \ill{}ssis eius $Z$ s. cuius quadrat. est $Z$ ss. si de eo auferatur quadrat. \uncertain{medius} $B$ plani cub. fit $Z$ ss. $-$ B plani cub. cuius latus est $L.\,Z$ ss. $-$ B plani cub. quod ablatum a semisse \uncertain{binomii} \struck{\ill{}} $Z$ s. dat $L\,V.\,Z$ s. $- L\,Z$ ss. $- Bppp$ \uncertain{extremum minus} ad eidem semissi additum dat maius extremum nempe $LVC.\,Z$ s. $+ L.\,Z$ ss. $- Bppp$ quæ iuncta faciunt extremorum summam æqualem ipsi $A$.}
\note{Note marginale d'une écriture plus petite et plus pâle, dans la marge de gauche, en regard du corollaire et de l'exemple qui suit ; elle justifie la règle : $Z$ solide deux fois est la somme de deux cubes extrêmes dont le produit est le cube de $B$ plan. Plusieurs mots de ses deux premières lignes sont sous des taches d'encre ; « semissis » n'est lu que par sa fin. Le mot devant « Z s. dat », en tête de ligne, est souligné et récrit au-dessus (« \ill{} ») ; « extremum minus » est écrit au bout de la ligne, en petit, et se lit mal. « Bppp » : $B$ plano-plano-plano. Un trait horizontal sépare cette note de la suivante.}

Si $1C - 81N$ æquetur $756$.

quoniam $378 + 351$ est $729$ cubus a latere $9$

et $378 - 351$ est $27$ cubus a latere $3$.

Ideo $9 + 3$ id est $12$ est $1N$ de qua quæritur
\note{L'exemple est juste : $378^2 - 27^3 = 123201 = 351^2$, et $12^3 - 81 \cdot 12 = 756$. Le « et » en tête de la deuxième ligne est surchargé.}

\marginal{Eadem sequuntur si \uncertain{negatus} \uncertain{suppositis} sic proponatur $\left.\begin{array}{l} Bp.\ \text{in } A\,3 \\ -\,Ac \end{array}\right\} Zs.$}
\note{Seconde note marginale, d'une encre plus noire et d'un trait plus appuyé, sous le trait qui clôt la première, en regard de l'exemple. Les deux mots lus avec doute sont écrits d'un trait rapide.}

\subsection*{De Canonica Æquationum transmutatione ut coefficientes subgraduales vel comparationum homogenea sint quæ præscribuntur. Cap VIII}
\note{Titre de trois lignes, centré, suivi de « Cap VIII », dont le dernier jambage se prolonge ; la lecture « VIIII » n'est pas exclue.}

Ad Impendia quoque Logistica confert æquationes ita præparare, ut coefficientes \struck{\ill{}} æquationum vel comparationum homogenea sint quæ præscribuntur quod libere \ill{} ex canonica transmutandi doctrina. Statuitur coefficiens unitas. Usus igitur impendij vel ex eo liquet quod potestates adfectæ sub unitate et gradu quocunque si modo radix est numeralis non aliter resolvuntur ac si essent puræ, neque enim necessario conturbat habenda alioqui (sed quæ \uncertain{debit} \uncertain{operis} \uncertain{salva} fit) parabolarum quarum quæque unitas est ratio. Et si radix non est numerus, ipsam radicem non esse numerum statim convincitur.
\marginal{Ex eo opere}
\marginal{\uncertain{Indutas}}
\marginal{ipsum \uncertain{$1N$} de qua quæritur}
\note{« quod libere \ill{} » : le mot en tête de ligne se lit « Aret », qui n'est pas latin. « numeralis » est écrit « numerãl » ; un « est » le précède, surchargé. « necessario » : le mot est corrigé en cours d'écriture. Dans « sed quæ debit operis salva fit », les mots « debit operis » sont soulignés, et la marge de gauche porte en regard « Ex eo opere », souligné, qui les remplace sans doute ; la syllabe finale de « salva » est surchargée, et « fit » souligné. Au-dessous, dans la marge, un mot souligné seul, lu « Indutas » ou « Innutas », sans appel visible dans le texte. « ipsam radicem » est souligné, avec un signe d'insertion, et la marge porte en regard « ipsum / X de qua quæritur » ; le signe lu $1N$ a la forme d'un trait oblique suivi d'un X. Une petite croix est tracée dans la marge au-dessous. La page s'arrête sur « convincitur » ; le reste du feuillet est blanc.}

\page{123}
\note{Verso du feuillet numéroté 57 (vue 122) : aucun numéro en tête, et la page poursuit le chapitre VIII commencé au bas de la vue 122 par ses exemples. Le feuillet est plus petit que le fond ; à droite, le bord d'un autre feuillet, avec quelques traits et signes de ses marges. Dans la marge de gauche, l'écriture du recto transparaît en miroir.}

$1C + 1N$ æquatur $10$, quoniam proximus cubus est $8$ cuius radix est $2$ quæ ducta in unitatem et admixta $8$ facit $10$ ideo radix quæsita est $2$. \uncertain{Atque} $1C - 1N$ æquetur $24$ quoniam proximus minor cubus est $8$ cuius radix est $2$ et adscita unitate $3$ sub qua et unitate facto solido cum multatur cubus ex $3$ relinquitur $24$ ideo radix quæsita est $3$
\marginal{æque.}
\note{« Atque » est souligné, et « æque. » écrit en regard dans la marge de gauche, qui le remplace sans doute. Un petit signe est tracé au-dessus du « 8 » de « minor cubus est 8 ». Les deux exemples sont justes : $2^3 + 2 = 10$, et $3^3 - 3 = 24$.}

Proponatur autem $1C + 1N$ æquari $9$ quoniam proximus cubus est $8$ cuius radix est $2$ quæ ducta in unitatem et admixta $8$ facit $10$ non etiam $9$, $1N$ est radix irrationalis

Atque $1C - 1N$ æquetur $25$ quoniam proximus minor cubus est $8$ cuius radix est \struck{\ill{}} $2$ et adscita unitate $3$ sub qua et unitate facto solido cum multatur cubus ex tribus relinquitur $24$ non etiam $25$ ideo $1N$ est irrationalis.

Oportet autem ad huiusmodi instituendæ transmutationis opus coefficientem quæ imperatur coefficienti æquationis propositæ esse congeneam et siquidem radix æquationis propositæ est eiusdem quoque generis concedetur esse. ut coefficiens propositæ æquationis ad coefficientem imperatam ita radix de qua quærebatur ad novam statuendam radicem.
\note{« opus » est surchargé d'une tache. « congeneam » : lecture des lettres, le mot étant écrit d'un trait (« congenerem » n'est pas exclu).}

Sin coefficiens gradus radicis generis communis concedetur esse. ut coefficiens propositæ æquationis ad coefficientem imperatam ita gradus æqualtus radicis de qua quærebatur ad gradum æqualtum novæ statuendæ radicis Et per resolutionem concessi analogismi exhibetur sub nova specie valor radicis quæsitæ et sic proposita æquatio dirigetur et ordinabitur nova.
\note{« æqualtus », « æqualtum » : ainsi écrit deux fois, en un mot (« æque altus », de même hauteur). La fin de « æquatio » est sous une tache.}

quod ut uno vel altero exemplo fiat apertius

Proponatur $A$ cubus plus $B$ in $A$ quad. æquari $Z$ solido

placeat autem æquationem ita transmutare ut affectio quidem maneat sub quadrato ipsaque adfirmata sed coefficiens sit $X$ non etiam $B$.

esto ut $B$ ad $X$ ita $A$ ad $E$.

ergo $\dfrac{B \text{ in } E}{X}$ erit $A$

quare secundum ea quæ proponuntur
\[ \left.\begin{array}{l} \dfrac{B\ \text{cubus in } E\ \text{cubum}}{X\ \text{cubo}} \\[2ex] +\,\dfrac{B\ \text{cubo in } E\ \text{quadrat.}}{X} \end{array}\right\}\ \text{æquabitur } Z\ \text{solido.} \]
\marginal{Xq.}
\note{Le second terme devrait avoir pour dénominateur $X$ quadrato ($B A^2$ avec $A = BE/X$) ; la ligne porte $X$, souligné, et la marge de gauche porte au bas, en regard, « Xq. », qui est sans doute la correction. Le bas du feuillet est blanc.}

\page{124}
\note{Comme à la vue 122, deux feuillets superposés. En haut, la même bande du feuillet numéroté « 59 » (« omnia applicentur ad $A + B$ » ; « Et oritur $A$ quadrato $- B$ in $A$ $\}$ fit $B$ quadratum bis »), transcrite à la vue 126. Posé dessus, un feuillet plus petit, numéroté « 58 » à l'encre en haut à droite, dont on transcrit la page. À gauche de la vue, le bord d'un autre feuillet, avec des fins de lignes qu'on ne transcrit pas. La page enchaîne sur l'égalité qui clôt la vue 123.}

Et omnibus per $X$ cubum ductis et $B$ cubum divisis
\[ \left.\begin{array}{l} E\ \text{cubus} \\ +\,X \text{ in } E\ \text{quadratum} \end{array}\right\}\ \text{æquabitur } \dfrac{X\ \text{cubo in } Z\ \text{solidum}}{B\ \text{cubo}} \]
Ipsum igitur factum est quod oportuit

Proponatur $1C + 20Q$ æquari $96{,}000$

$1C + 1Q$ æquabitur $12$ et fit $1N$ $2$

unde fit radix prima quæsita $40$.
\note{L'exemple est juste : avec $B = 20$ et $X = 1$, $96000/8000 = 12$, $2^3 + 2^2 = 12$, et $A = 20 \cdot 2 = 40$ ; $40^3 + 20 \cdot 40^2 = 96000$.}

\subsection*{Aliud.}

Proponatur $A$ cubus minus $B$ quad. in $A$ æquari $Z$ solido.

placeat autem æquationem ita transmutare ut affectio quidem maneat sub latere ipsaque negata sed coefficiens sit $X$ quadratum non etiam $B$ quadratum.

Esto ut $B$ quad. ad $X$ quad. ita $A$ quad. ad $E$ quad. et consequenter ut $B$ ad $X$ ita $A$ ad $E$

$\dfrac{B \text{ in } E}{X}$ erit $A$ quare secundum ea quæ proponuntur
\[ \left.\begin{array}{l} \dfrac{B\ \text{cubus in } E\ \text{cubum}}{X\ \text{cubo}} \\[2ex] -\,\dfrac{B\ \text{cubo in } E}{X} \end{array}\right\}\ \text{æquabitur } Z\ \text{solido} \]
Et omnibus per $X$ cubum ductis et $B$ cubum divisis
\[ \left.\begin{array}{l} E\ \text{cubus} \\ -\,X\ \text{quadrato in } E \end{array}\right\}\ \text{æquabitur } Z\ \text{solido} \]
\marginal{$\dfrac{Xc \text{ in } Zs.}{Bc}$}
\note{« Z solido » est souligné, et la marge de gauche porte en regard, souligné aussi, « Xc in Z s. / Bc » : la correction, puisque le second membre doit être multiplié par $X$ cube et divisé par $B$ cube, comme dans le premier exemple.}

Ipsum igitur factum est quod oportuit

Proponatur $1C - 144N$ æquari $10368$

$1C - 1N$ æquabitur $6$ et est $1N$ $2$

unde fit radix prima quæsita $24$.
\note{L'exemple est juste : avec $B = 12$ et $X = 1$, $10368/1728 = 6$, $2^3 - 2 = 6$, et $A = 12 \cdot 2 = 24$.}

Neque vero opus aliter fit in præscriptis comparationum homogeneis concedetur nimirum esse ut magnitudo æqualis potestati quæ proponitur adfectæ ad magnitudinem præscriptam æqualtum et homogeneam ita potestas radicis de qua quærebatur ad potestatem novæ statuendæ radicis \struck{ex} per resolutionem concessi analogismi exhibetur sub nova specie valor radicis quæsitæ et proposita æqualitas dirigetur et ordinabitur nova. Exempli causa
\note{« opus » est surchargé ; « æqualtum » est écrit ainsi, au masculin, auprès de « præscriptam » et « homogeneam » ; voir la vue 123 ; « analogismi » est sous une tache. Le mot barré devant « per » est bref.}

\[ \text{Proponatur } \left.\begin{array}{l} A\ \text{cubus} \\ \text{plus } B\ \text{plano in } A. \end{array}\right\}\ \text{æquari } Z\ \text{cubo} \]
\note{La page s'arrête sur cet énoncé ; le reste du feuillet est blanc.}

\page{125}
\note{Verso du feuillet numéroté 58 (vue 124) : aucun numéro en tête ; le feuillet, plus petit que le fond, laisse voir au-dessus le bord blanc d'un autre, et à droite un onglet. La page poursuit l'exemple posé au bas de la vue 124 (« Proponatur $A$ cubus plus $B$ plano in $A$ æquari $Z$ cubo »).}

placeat autem æquationem ita transmutare ut potestas adfirmate adfecta sub radice et coefficiente $B$ plano comparetur $D$ cubo

Concedetur esse ut $Z$ ad $D$ ita $A$ ad $E$

ergo $\dfrac{Z \text{ in } E}{D}$ erit $A$
\[ \text{quare } \left.\begin{array}{l} \dfrac{Z\ \text{cubus in } E\ \text{cubum}}{D\ \text{cubo}} \\[2ex] +\,\dfrac{B\ \text{plano in } Z \text{ in } E}{D} \end{array}\right\}\ \text{æquabitur } Z\ \text{cubo} \]
Et omnibus in $D$ cubum ductis et per $Z$ cubum divisis
\[ \left.\begin{array}{l} E\ \text{cubus} \\ +\,\dfrac{B\ \text{plano in } D\ \text{quad. in } E}{Z\ \text{quad.}} \end{array}\right\}\ \text{æquabitur } D\ \text{cubo} \]
Ipsum igitur factum est quod oportuit
\note{« Concedetur » est écrit « Concreditor » ou « Concedetor » ; lecture par le sens, comme aux vues 123 et 124. Le « B » de « B plano in D quad. » est surchargé.}

Proponatur $1C + 860N$ æquari $1{,}728$

$1C + 215N$ æquabitur $216$ et est $1N$ $1$

unde fit radix prima quæsita $2$.
\note{L'exemple est juste : avec $Z = 12$ et $D = 6$, $860 \cdot 36/144 = 215$, $1 + 215 = 216$, et $A = 12 \cdot 1/6 = 2$ ; $8 + 1720 = 1728$.}

\subsection*{Anomala æquationum aliquot cubicarum ad quadraticas aut etiam simpliciores reductio. Cap IX.}
\marginal{IX}
\note{« aut etiam » est souligné. « Cap IX » est souligné, et « IX », souligné aussi, est répété dans la marge de gauche.}

Præparandarum igitur æquationum solennes modi ita se habent. de irregularibus autem non statuitur præcepta quoniam anomalia \struck{illa} non magis finita est quam artificis in indagando vis et solertia. ad excitandam tamen eam vim et solertiam opportunum est singularia aliquot æquationum constitutiva et reductiva theoremata adnotasse, insignem aliquam emphasim vel elegantiam præseferentia, qualia iam sequuntur
\note{« statuitur » est écrit tel ; « præcepta » suit. Le mot barré devant « non magis » est surchargé ; « illa » est une lecture. Une petite croix est tracée dans la marge de gauche en regard de cet alinéa.}

\subsection*{prop. 1.}

\[ \text{Si } \left.\begin{array}{l} A\ \text{cubus} \\ -\,B\ \text{quadrato bis in } A \end{array}\right\}\ \text{æquetur } B\ \text{cubo} \]
\[ \left.\begin{array}{l} A\ \text{quadratum} \\ -\,B \text{ in } A \end{array}\right\}\ \text{æquabitur } B\ \text{quadrato} \]
Ex iis \struck{\ill{}}enim quæ proponuntur manifestum fit per antithesin $A$ cubum æquari $B$ cubo $+$ $B$ quadrato bis in $A$. et addendo utrique parti $B$ cubum
\[ \left.\begin{array}{l} A\ \text{cubum} \\ +\,B\ \text{cubum} \end{array}\right\}\ \text{æquari } \begin{array}{l} B\ \text{cubo bis} \\ +\,B\ \text{quadrato bis in } A \end{array} \]
\note{Le début de « enim » est noirci sous une rature. La page s'arrête sur cette égalité, et le reste du feuillet est blanc ; la suite (« omnia applicentur ad $A + B$ ») est en tête du feuillet numéroté 59, dont le haut paraît aux vues 122 et 124.}

\page{126}
\note{Le feuillet numéroté « 59 » à l'encre en haut à droite, dont le haut paraît déjà aux vues 122 et 124, est ici vu en entier. Il est plus court que les autres : sous son bord inférieur paraît le bas d'un feuillet plus grand, avec trois lignes d'une encre plus pâle (« quoniam enim $A$ quadratum $+ B$ in $A$ bis $\}$ æquatur $B$ quadrato bis $+ D$ plano »), qui ne sont pas de cette page et ne sont pas transcrites ici. La page continue la proposition 1 du chapitre IX laissée au bas de la vue 125. Dans les marges, l'écriture de l'autre face transparaît en miroir.}

omnia applicentur ad $A + B$
\[ \text{Et oritur } \left.\begin{array}{l} A\ \text{quadrat.} \\ -\,B \text{ in } A \\ +\,B\ \text{quadrato} \end{array}\right\}\ \text{fit } B\ \text{quadratum bis} \]
et consequenter ablato utrinque $B$ quadrato
\[ \left.\begin{array}{l} A\ \text{quadratum} \\ -\,B \text{ in } A \end{array}\right\}\ \text{æquabitur } B\ \text{quadrato} \]
Si $1C - 18N$ æquetur $27$ igitur $1Q - 3N$ æquabitur $9$.
\note{Sous « fit », un petit signe (« b » ?) sans appel visible. L'exemple suit la règle avec $B = 3$.}

\subsection*{prop. II.}

\[ \text{Si } \left.\begin{array}{l} B\ \text{quadrat. bis in } A \\ -\,A\ \text{cub.} \end{array}\right\}\ \text{æquetur } B\ \text{cubo} \]
\[ \left.\begin{array}{l} A\ \text{quadrat.} \\ +\,B \text{ in } A \end{array}\right\}\ \text{æquabitur } B\ \text{quadrato} \]
Ex iis enim quæ proponuntur manifestum \struck{est} fit per antithesin $A$ cubum æquari $B$ quad. bis in $A - B$ cubo et auferendo utrique parti $B$ cubum
\[ \left.\begin{array}{l} A\ \text{cubum} \\ -\,B\ \text{cubo} \end{array}\right\}\ \text{æquari } B\ \text{quadrato bis in } A \;-\; B\ \text{cubo bis} \]
omnia applicentur ad $A - B$
\[ \text{Et oritur } \left.\begin{array}{l} A\ \text{quadratum} \\ +\,B \text{ in } A \\ +\,B\ \text{quadrato} \end{array}\right\}\ \text{fit } B\ \text{quadratum} \]
\marginal{bis}
\note{« B quadratum » est souligné, et « bis », souligné, est écrit en regard dans la marge de gauche : $B$ quadratum bis, comme à la proposition 1.}
et consequenter ablato utrinque $B$ quadrato
\[ \left.\begin{array}{l} A\ \text{quadratum} \\ +\,B \text{ in } A \end{array}\right\}\ \text{æquabitur } B\ \text{quadrato} \]
Si $18N - 1C$ æquatur $27$ Igitur $1Q - 3N$ æquabitur $9$.
\note{La règle qui précède donne $1Q + 3N$ ; la ligne porte nettement « $-$ ».}

\subsection*{prop. III}

Si $B$ quadrat. in $A$ ter minus $A$ cubo, æquetur $B$ cubo bis

$B$ dupla est ipsius $A$ de qua quæritur
\marginal{Si $\left.\begin{array}{l} A\ \text{cub.} \\ -\,Bq\ \text{ter in } A \end{array}\right.$ ; ipsa}
\note{Dans la marge de gauche, en regard, deux corrections soulignées : « Si A cub. $-$ Bq ter in A », qui renverse les deux termes de l'énoncé (l'exemple qui suit est bien $1C - 12N$), et « ipsa », en regard de « ipsius », souligné dans la ligne. Le « A » de « A cubo » est surchargé d'une tache.}

quoniam enim $B$ est dupla ipsius $A$ de qua quæritur

ergo ex iis quæ proponuntur $B$ cubus $8 - B$ quadrato in $B$ $6$ æquabitur $B$ cubo bis quod quidem ita se habet
\note{Un petit signe de renvoi en forme de croix est tracé à gauche de « quoniam » ; il répond au même signe dans la note marginale qui suit.}

$1C - 12N$ æquatur $16$ fit $1N$ $4$.

$1C - 6N$ æquat. $L\,32$ fit $1N$ $L\,2$.
\note{Le premier exemple est juste ($B = 2$, $A = 4$). Au second, $B$ est $L\,2$ et la racine de $x^3 - 6x = \sqrt{32}$ est $2\sqrt{2} = L\,8$ ; la ligne porte « $L\,2$ », qui est $B$. Un point est posé au-dessus du « 6 ».}

\marginal{prop. III. Si $\left.\begin{array}{l} A\ \text{cubus} \\ -\,Bq.\ \text{ter in } A \end{array}\right\}$ æq. $Bc\,2$ $B$ dupla est ipsius $A$ ; quoniam enim $B\,2 \} A$. ergo $\left.\begin{array}{l} Bc\,8. \\ -\,Bq \text{ in } B\,6. \end{array}\right\} Bc$ bis ; $1C - 12N$ æquat. $16$ fit $1N$ $4$. ; $1C - 3N$ æquatur $2$ fit $1N$ $2$.}
\note{Note marginale de la main des corrections, d'une encre plus noire, sous un trait qui la sépare des précédentes : la proposition 3 récrite. Elle pose $A$ double de $B$ (« $B\,2 \} A$ »), ce qui est juste ($8B^3 - 6B^3 = 2B^3$), bien que la phrase garde « B dupla est ipsius A », dont le dernier mot est surchargé. Les deux exemples chiffrés sont écrits au bas de la marge, sur toute la largeur du pied du feuillet, soulignés d'un trait courbe ; le second, $1C - 3N = 2$, $1N = 2$, est juste.}

\marginal{propos. IIII ; Si $\left.\begin{array}{l} Bq\ 3 \text{ in } A \\ -\,Ac \end{array}\right\}$ æquetur $Bc$ bis $B$ \struck{subdupla est ipsius} \add{æquatur ipsi} $A$ de qua quæritur quoniam enim $B$ \struck{dupla est ipsius} \add{æquatur ipsi} $A$ \struck{$Bc\,8 - Bq\,3 \text{ in } B\,2 \} Bc$ bis} $\left.\begin{array}{l} Bq \text{ in } B\,3 \\ -\,Bc. \end{array}\right\} Bc\,2$. ; $12N - 1C$ æquatur $16$ fit $1N$ $2$ ; et vel $6N - 1C$ æquat. $L\,32$ fit $1N$ $L\,2$.}
\note{Au pied du feuillet, sur deux lignes serrées qui courent d'un bord à l'autre, de la même main noire : une proposition 4. Les mots barrés sont récrits au-dessus de la ligne (« æquatur ipsi », deux fois). Le calcul barré est illisible en partie ; la lecture donnée est douteuse. Les exemples sont justes : $3B^3 - B^3 = 2B^3$ pour $A = B$, $24 - 8 = 16$, et $6\sqrt{2} - 2\sqrt{2} = \sqrt{32}$. C'est de cette proposition que le second exemple de la proposition 3 (« fit 1N L2 ») semble repris.}

\page{127}
\note{Verso du feuillet numéroté 59 (vue 126) : aucun numéro en tête ; en haut à gauche, un « 20 » vu en miroir, qui est une transparence. Le feuillet est plus court que le fond : sous son bord inférieur paraît le bas d'un autre feuillet, avec la dernière égalité de la vue 125 (« $A$ cubum $+ B$ cubum \ldots{} $+ B$ quadrato bis in $A$ »), qui n'est pas de cette page. À droite, le bord d'un autre feuillet, avec des bouts de mots. La page poursuit les propositions du chapitre IX ; la proposition IIII n'est écrite qu'au pied du recto (vue 126), de la main des corrections.}

\subsection*{prop. V.}

Si $A$ cubus minus $B$ in $A$ quadratum plus $D$ plano in $A$ æquetur $B$ in $D$ plano, ipsa $B$ est $A$ de qua quæritur

quoniam enim ipsa $B$ est $A$ de qua quæritur ergo ex iis quæ proponuntur $B$ cubus minus \uncertain{$B$} in $B$ quadratum plus $D$ plano in $B$ æquabitur $B$ in $D$ plano
\marginal{$B$}
\note{La lettre devant « in B quadratum » est noyée sous une tache ; la marge de gauche porte en regard un « B » souligné, qui la rétablit.}

quod quidem manifeste ita se habet

$1C - 4Q + 5N$ æquatur $20$ facto ex $4$ in $5$ ergo $1N$ fiet $4$.
\note{L'exemple est juste : $64 - 64 + 20 = 20$.}

\subsection*{prop. VI}

Si $A$ cubus plus $B$ in $A$ quadratum minus $D$ quadrato in $A$ æquetur $B$ in $D$ quadratum. ipsa $D$ est $A$ de qua quæritur

quoniam enim ipsa $D$ est $A$ de qua quæritur ergo ex iis quæ proponuntur $B$ cubus, plus $B$ in $D$ quadrat. minus $D$ quadrato in $D$ æquabitur $B$ in $D$ quadratum.
\marginal{$D$}
\note{« B cubus » est souligné, et la marge de gauche porte en regard un « D » souligné : il faut « D cubus », puisque $D$ est mis pour $A$.}

quod quidem manifeste ita se habet

$1C + 5Q - 4N$ æquatur $20$ facto ex $5$ in $4$ ergo $1N$ fit $L\,4$.
\note{L'exemple est juste : $D$ quadratum est $4$, et $8 + 20 - 8 = 20$ pour $1N = L\,4 = 2$.}

\subsection*{prop. VII}

Si $B$ in $A$ quadrat. plus $D$ quadrato in $A$ minus $A$ cubo æquetur $D$ quadrato in $B$ ipsa $B$ vel $D$ est $A$ de qua quæritur

quoniam enim ipsa $B$ est $A$ de qua quæritur ergo ex iis quæ proponuntur $B$ cubus plus $D$ quadrato in $B$ minus $B$ cubo æquabitur $B$ in $D$ quadratum quod quidem ita manifeste se habet

rursus quoniam ipsa $D$ est $A$ de qua quæritur ergo ex iis quæ proponuntur $B$ in $D$ quadratum plus $D$ cubo minus $D$ cubo æquabitur $D$ quadrato in $B$ quod quidem ita quoque manifeste se habet

$6Q + 4N - 1C$ æquatur $24$ fit $1N$ $6$ vel $2$
\note{L'exemple est juste pour les deux racines : $216 + 24 - 216 = 24$ et $24 + 8 - 8 = 24$.}

\subsection*{prop. VIII}

\[ \text{Si } \left.\begin{array}{l} D \text{ in } A\ \text{quad.} \\ \text{plus } B \text{ in } D \text{ in } A \\ \text{minus } A\ \text{cubo} \end{array}\right\}\ \text{æquetur } B\ \text{cubo} \]
\[ \left.\begin{array}{l} D \text{ in } A \\ =\,B \text{ in } A \\ -\,A\ \text{quadrato} \end{array}\right\}\ \text{æquabitur } B\ \text{quadrato} \]
\note{Devant « B in A », deux traits parallèles, « = », là où les autres lignes ont un seul trait ; le signe est laissé tel. Le calcul demande $+B$ in $A$ : $A^2 - (B + D)A + B^2$ est le quotient par $A + B$. Une croix dans la marge de gauche, en regard. La page s'arrête ici ; la proposition n'a ni démonstration ni exemple sur ce feuillet.}

\page{128}
\note{Feuillet numéroté « 60 » à l'encre en haut à droite. Comme le 59, il est plus court que le fond : sous son bord inférieur reparaissent les trois lignes pâles déjà vues sous le 59 à la vue 126 (« quoniam enim $A$ quadratum $+ B$ in $A$ bis \ldots{} $+ D$ plano »), qui ne sont pas de cette page. À gauche, le bord du feuillet précédent, avec des fins de lignes ; en haut à gauche, de l'écriture vue en miroir, qui est une transparence. La page continue la proposition VIII laissée au bas de la vue 127.}

Ex iis enim quæ proponuntur manifestum fit per antithesin
\[ \left.\begin{array}{l} B\ \text{cubum} \\ +\,A\ \text{cubo} \end{array}\right\}\ \text{æquari } \begin{array}{l} D \text{ in } A\ \text{quad.} \\ +\,D \text{ in } A \text{ in } B \end{array} \]
utraque pars æqualitatis applicetur ad $A + B$
\[ \text{ergo } \left.\begin{array}{l} A\ \text{quad.} \\ -\,B \text{ in } A \\ +\,B\ \text{quad.} \end{array}\right\}\ \text{æquabitur } D \text{ in } A \]
\[ \text{Et per antithesin } \left.\begin{array}{l} D \text{ in } A \\ +\,B \text{ in } A \\ -\,A\ \text{quadrato} \end{array}\right\}\ \text{æquabitur } B\ \text{quadrato} \]
\note{Ici le signe devant « B in A » est un « + » net ; voir la vue 127. Une tache d'encre sur « manifestum ».}

Si $1Q + 20N - 1C$ æquetur $8$ quia latus cubi $8$ ductum in $10$ facit $20$ igitur $12N - 1Q$ æquabitur $4$ et fit $1N$ $6 - L\,32$ vel $6 + L\,32$.
\marginal{10}
\note{« 1Q » est souligné, et « 10 », souligné, est écrit en regard dans la marge de gauche : il faut $10Q$, puisque $D$ vaut $10$. Avec $B = 2$ et $D = 10$, l'exemple est juste : $A^2 - 12A + 4 = 0$ donne $6 \pm \sqrt{32}$.}

\subsection*{prop. IX}

\[ \text{Si } \left.\begin{array}{l} A\ \text{cub.} \\ -\,D \text{ in } A\ \text{quad.} \\ +\,D \text{ in } B \text{ in } A \end{array}\right\}\ \text{æquetur } B\ \text{cubo} \]
\[ \left.\begin{array}{l} D \text{ in } A \\ -\,B \text{ in } A \\ -\,A\ \text{quadrato} \end{array}\right\}\ \text{æquabitur } B\ \text{quadrato} \]
Ex iis enim quæ proponuntur manifestum fit per antithesin cum $A$ intelligitur maior quam $B$
\[ \left.\begin{array}{l} A\ \text{cubum} \\ -\,B\ \text{cubo} \end{array}\right\}\ \text{æquari } \begin{array}{l} D \text{ in } A\ \text{quadrat.} \\ -\,D \text{ in } A \text{ in } B \end{array} \]
utraque pars æqualitatis applicetur ad $A - B$
\[ \text{Igitur } \left.\begin{array}{l} A\ \text{quadratum} \\ +\,B \text{ in } A \\ +\,B\ \text{quadrato} \end{array}\right\}\ \text{æquabitur } D \text{ in } A \]
\[ \text{Et per antithesin } \left.\begin{array}{l} D \text{ in } A \\ -\,B \text{ in } A \\ -\,A\ \text{quadrato} \end{array}\right\}\ \text{æquabitur } B\ \text{quadrato} \]
Atque cum $B$ intelligitur maior quam $A$
\[ \left.\begin{array}{l} B\ \text{cubus} \\ -\,A\ \text{cubo} \end{array}\right\}\ \text{æquabitur } \begin{array}{l} D \text{ in } A \text{ in } B \\ -\,D \text{ in } A\ \text{quadrat.} \end{array} \]
utraque pars æqualitatis applicetur ad $B - A$
\[ \left.\begin{array}{l} A\ \text{quadrat.} \\ +\,B \text{ in } A \\ +\,B\ \text{quadrato} \end{array}\right\}\ \text{æquabitur } D \text{ in } A \text{ ut ante} \]
\note{Le long trait qui traverse « B cubus » part du jambage de « Atque » à la ligne au-dessus ; ce n'est pas une rature. La page s'arrête ici ; le bas du feuillet est blanc.}

\page{129}
\note{Verso du feuillet numéroté 60 (vue 128) : aucun numéro en tête ; en haut à gauche, « 60 » vu en miroir, transparence du recto. Sous le bord inférieur du feuillet reparaît, sur le feuillet de dessous, la dernière égalité de la vue 125, qui n'est pas de cette page. La page s'ouvre sur l'exemple de la proposition IX (vue 128).}

Si $1C + 10Q + 20N$ æquetur $8$ quia $LC\,8$ ductum in $10$ facit $20$. igitur \ill{}$8N - 1Q$ æquabitur $4$ et fit $1N$ $4 - L\,12$ vel $4 + L\,12$.
\note{« 10Q » est souligné. Selon la proposition IX, l'exemple devrait être $1C - 10Q + 20N$ ; le signe lu devant « 10Q » est un « + ». Une tache couvre le début de « 8N ». Avec $B = 2$ et $D = 10$, la suite est juste : $A^2 - 8A + 4 = 0$ donne $4 \pm \sqrt{12}$.}

\subsection*{prop. \struck{10} X.}

\[ \text{Si } \left.\begin{array}{l} A\ \text{cubus} \\ -\,B\ \text{plano ter in } A \end{array}\right\}\ \text{æquetur } L\,B\ \text{plano-plano-plano bis} \]
\marginal{1}
\[ \left.\begin{array}{l} \dfrac{\begin{array}{l} L\,B\ \text{plani}\ 3 \\ \pm\,L\,B\ \text{plani}\ 1 \end{array}}{L\,2} \end{array}\right\}\ \text{fit } A \text{ de qua quæritur} \]
\note{Le numéro de la proposition est d'abord écrit « 10 », barré, puis « X ». « L B plani 3 » : le chiffre, écrit comme « le », est lu 3 d'après le calcul qui suit ($L\,27$ pour son cube). Devant « L B plani 1 », un signe « $\pm$ » (un « + » sur deux traits) ; en regard, dans la marge de gauche, un signe noirci d'encre. Plus haut, un « 1 » souligné dans la marge, en regard de l'énoncé, sans appel visible ; « plano » de « plano-plano-plano » est souligné. Le premier membre de l'énoncé est précédé d'un trait (« $-$ Si »).}

\[ \text{quoniam enim } \left.\dfrac{\begin{array}{l} L\,B\ \text{plani}\ 3 \\ +\,L\,B\ \text{plani}\ 1 \end{array}}{L\,2}\right\}\ \text{est ipsa } A \]
\marginal{de qua quæritur}
\note{« A » est souligné, et « de qua quæritur » est écrit en regard dans la marge de gauche, qui complète la phrase.}

\[ \text{ideo ex iis quæ proponuntur } \left.\begin{array}{l} \dfrac{\begin{array}{l} L\,B\ \text{plano-plano-plani}\ 27 \\ +\,L\,B\ \text{plano-plano-plani}\ 81 \\ +\,L\,B\ \text{plano-plano-plani}\ 27 \\ +\,L\,B\ \text{plano-plano-plani}\ 1 \end{array}}{L\,8} \\[4ex] \dfrac{\begin{array}{l} -\,L\,B\ \text{plano-plano-plani}\ 27 \\ -\,L\,B\ \text{plano-plano-plani}\ 9 \end{array}}{L\,2} \end{array}\right\}\ \text{æquat. } L\,B\ \dfrac{\text{plano-plano-plani}}{.1.} \]
\marginal{2}
\note{Le second membre est souligné, avec « .1. » sous le trait, et un « 2 » souligné est écrit en regard dans la marge de gauche : c'est la correction, le calcul donnant bien $\sqrt{2B^3}$ (le cube de $(\sqrt{3} + 1)/\sqrt{2}$ moins trois fois ce nombre vaut $\sqrt{2}$). Une tache couvre la lettre après « L ». Plus haut dans la marge, un signe « $\frac{0}{X}$ », peut-être un renvoi.}

quod quidem ita se habet subducendo in prima æqualitatis parte ab æqualibus æqualia

$1C - 6N$ æquatur $4$ igitur $L\,3 - 1$ fit $1N$.
\marginal{$\pm$}
\note{« $L\,3 -$ » est souligné, et la marge de gauche porte en regard le même signe « $\pm$ » qu'au-dessus. La racine positive de $x^3 - 6x = 4$ est $\sqrt{3} + 1$ ; $\sqrt{3} - 1$ est celle de la proposition suivante.}

\subsection*{prop. XI.}

\[ \text{Si } \left.\begin{array}{l} B\ \text{plano ter in } A \\ -\,A\ \text{cubo} \end{array}\right\}\ \text{æquetur } L\,B\ \text{plano-plano-plano}\ 2 \]
\marginal{1}
\[ \left.\dfrac{\begin{array}{l} L\,B\ \text{plani}\ 3 \\ -\,L\,B\ \text{plani}\ 1 \end{array}}{L\,2}\right\}\ \text{fit } A \text{ minor de qua quæritur} \]
ut apparet insequendo vestigia antecedentis demonstrationis

$6N - 1C$ æquatur $4$

igitur \struck{\ill{}} $L\,3 - 1$ fit $1N$ \struck{eaque} minor altera est $2$.
\note{Un « 1 » souligné dans la marge, en regard de l'énoncé, comme à la proposition X ; « plano » est souligné. Devant « L3 », une lettre surchargée ; « eaque » est barré d'un trait, et la lecture du mot barré est douteuse. L'exemple est juste : $6x - x^3 = 4$ a pour racines $2$ et $\sqrt{3} - 1$.}

\marginal{quad. \struck{cubus} $\left.\begin{array}{l} L\,Bp.\ 3 \\ +\,L\,Bpl.\ 1 \end{array}\right.$ est $\dfrac{\begin{array}{l} Bpl.\ 3 \\ +\,L\,Bpp.\ 12 \\ +\,Bp. \end{array}}{2}$ et cubus $\dfrac{\begin{array}{l} L\,Bppp.\ 27 \\ +\,L\,Bppp.\ 36 \\ +\,L\,Bppp.\ 3 \\ +\,L\,Bppp.\ \ill{} \\ +\,L\,Bppp.\ 12 \\ +\,L\,Bppp.\ 1 \end{array}}{L\,8}$ ; $L\,81$, $L\,27$, $L\,16$ ; \uncertain{solum} ablatis $\left\{\dfrac{\begin{array}{l} L\,Bppp.\ 27 \\ -\,L\,Bppp.\ 9 \end{array}}{L\,2}\right.$ est autem \struck{\ill{}} $L\,8$. duplum ipsius $L\,2$. ergo ablatis æqualibus remanet $\dfrac{L\,Bppp.\ 16}{L\,8}$ id est $L\,Bppp.\ 2$.}
\note{Au pied de la page, sous un trait tiré sur toute la largeur et partant d'un signe « $\frac{0}{X}$ » dans la marge (le même qu'en regard de la vérification de la proposition X), un calcul de vérification qui court d'un bord à l'autre : le carré, puis le cube, de $(\sqrt{3B} + \sqrt{B})/\sqrt{2}$. Des traits croisés réunissent les termes du cube deux à deux et mènent à « $L\,81$ », « $L\,27$ » et « $L\,16$ » ; le quatrième terme du cube est sous une tache. Une accolade descend de la dernière ligne vers la gauche.}

\page{130}
\note{Feuillet numéroté « 61 » à l'encre en haut à droite. Au-dessus, sur une bande, le haut d'un autre feuillet, numéroté « 63 », avec trois lignes (« quod quidem ita se habet » ; « $30Q - 156N - 1C$ æquatur $440$ fit $1N$ $10$ » ; « prop. VI ») qui ne sont pas de cette page. Les trois dernières lignes de celle-ci (« quoniam enim $A$ quadratum $+ B$ in $A$ bis \ldots{} $+ D$ plano ») sont celles qu'on voyait sous le bord des feuillets 59 et 60 aux vues 126 et 128. Le feuillet 61 commence un nouveau chapitre ; la proposition IX de la vue 128 et les propositions X et XI de la vue 129 n'ont pas de suite ici.}

\subsection*{Similium reductionum Continuatio Cap. XX.}
\marginal{X.}
\note{« XX » est souligné, et « X. », souligné, est écrit en regard dans la marge de gauche : chapitre X, à la suite du chapitre IX de la vue 125.}

\subsection*{prop. 1}

\[ \text{Si } \left.\begin{array}{l} A\ \text{cubus} \\ +\,B \text{ in } A\ \text{quad. ter} \\ +\,D\ \text{plano in } A \end{array}\right\}\ \text{æquetur } \begin{array}{l} B\ \text{cubo bis} \\ -\,D\ \text{plano in } B \end{array} \]
\[ \left.\begin{array}{l} A\ \text{quadrat.} \\ +\,B \text{ in } A\ \text{bis} \end{array}\right\}\ \text{æquatur } \begin{array}{l} B\ \text{quadrat. bis} \\ -\,D\ \text{plano} \end{array} \]
\note{Dans l'énoncé, les deux « + » de la colonne de gauche sont tracés d'un long trait qui traverse le début de la ligne ; ce ne sont pas des ratures. Une tache sur « quadrat. » de « B quadrat. bis ».}

\[ \text{quoniam enim } \left.\begin{array}{l} A\ \text{quadratum} \\ +\,B \text{ in } A\ \text{bis} \end{array}\right\}\ \text{æquat. } B\ \text{quadrato bis} \]
\marginal{$-\,D$ plano \struck{in $A$}}
\note{À droite de « B quadrato bis », un trait et un petit signe (« x° ») ; la marge de gauche porte en regard « $- D$ plano », soulignée, suivie de « in A » barré : c'est le terme omis, $B$ quadrato bis $- D$ plano.}

ductis igitur omnibus in $A$
\[ \left.\begin{array}{l} A\ \text{cubus} \\ =\,B \text{ in } A\ \text{quadrat. bis} \end{array}\right\}\ \text{æquabitur } \begin{array}{l} B\ \text{quad. in } A\ 2 \\ -\,D\ \text{plano in } A \end{array} \]
\marginal{$+$}
\note{Le signe devant « B in A quadrat. bis » est un double trait, surchargé ; la marge de gauche porte en regard un « + », qui est le signe juste.}

et iisdem ductis in $B$
\[ \left.\begin{array}{l} B \text{ in } A\ \text{quadrat.} \\ +\,B\ \text{quadrato in } A\ \text{bis} \end{array}\right\}\ \text{æquabitur } \begin{array}{l} B\ \text{cubo bis} \\ -\,D\ \text{plano in } B \end{array} \]
Iungantur ducta æqualia æqualibus
\[ \left.\begin{array}{l} A\ \text{cubus} \\ =\,B \text{ in } A\ \text{quadrat. ter} \\ =\,B\ \text{quad. in } A\ \text{bis} \end{array}\right\}\ \text{æquabitur } \begin{array}{l} +\,B\ \text{quadrat. in } A\ \text{bis} \\ =\,D\ \text{plano in } A \\ +\,B\ \text{cubo bis} \\ -\,D\ \text{plano in } B \end{array} \]
\marginal{$+D$ ; $+$ ; $-$ ; $+$}
\note{Plusieurs signes de cette égalité sont des doubles traits, repris ou surchargés, et des traits partent de la marge de gauche vers eux ; la marge porte en regard « $+D$ » souligné, « $+$ », « $-$ », et plus bas un « $+$ » souligné. Le « + » de « $+$ B quadrat. in A bis » est écrit sur un autre signe. Tout est laissé comme il est écrit.}

et deleta utrinque adfectione $B$ quadrati in $A$ bis et ad æqualitatis ordinationem translata per antithesin
\marginal{\struck{\ill{} facta} lata per Antithesin $D$ plani in $A$ adfectione.}
\note{Sous « translata », un « A » écrit dans l'interligne. La note marginale, dont les deux premiers mots sont barrés, complète la phrase : l'affection de $D$ plan en $A$ passe de l'autre côté par antithèse.}
\[ \left.\begin{array}{l} A\ \text{cubus} \\ =\,B \text{ in } A\ \text{quadrat. ter} \\ +\,D\ \text{plano in } A \end{array}\right\}\ \text{æquabitur } \begin{array}{l} B\ \text{cubo bis} \\ -\,D\ \text{plano in } B. \end{array} \]
\note{Devant « B in A quadrat. ter », un double trait, avec un trait venu de la marge ; en regard, dans la marge, un signe noirci d'encre.}

quod quidem ita se habet

$1C + 30Q + 44N$ æquatur $1560$

Igitur $1Q + 20N$ æquabitur $756$ et fit $1N$ $6$.
\marginal{156}
\note{« 756 » est souligné, et « 156 », souligné, est écrit en regard dans la marge de gauche : c'est la correction. Avec $B = 10$ et $D = 44$ : $2000 - 440 = 1560$, $200 - 44 = 156$, et $36 + 120 = 156$ ; $216 + 1080 + 264 = 1560$.}

\subsection*{prop. II}

\[ \text{Si } \left.\begin{array}{l} A\ \text{cubus} \\ +\,B \text{ in } A\ \text{quadrat. ter} \\ -\,D\ \text{plano in } A \end{array}\right\}\ \text{æquetur } \begin{array}{l} B\ \text{cubo bis} \\ +\,D\ \text{plano in } B \end{array} \]
\[ \left.\begin{array}{l} A\ \text{quadratum} \\ +\,B \text{ in } A\ \text{bis} \end{array}\right\}\ \text{æquabitur } \begin{array}{l} B\ \text{quadrato bis} \\ +\,D\ \text{plano} \end{array} \]
\[ \text{quoniam enim } \left.\begin{array}{l} A\ \text{quadratum} \\ +\,B \text{ in } A\ \text{bis} \end{array}\right\}\ \text{æquatur } \begin{array}{l} B\ \text{quadrato bis} \\ +\,D\ \text{plano.} \end{array} \]
\note{La page s’arrête ici ; la démonstration continue en tête de la vue 131, verso du feuillet.}

\page{131}
\note{Verso du feuillet numéroté 61 (vue 130) : aucun numéro en tête ; en haut à gauche, « 61 » vu en miroir, transparence du recto. Le feuillet est plus petit que la vue ; à droite, le bord d'un autre feuillet. La page continue la démonstration de la proposition II du chapitre X, laissée au bas de la vue 130. Dans la marge de gauche, plusieurs signes de correction faits d'une croix posée sur un double trait (rendus ici « $\pm$ »), en regard des signes du texte qu'ils corrigent.}

ductis igitur omnibus in $A$
\[ \left.\begin{array}{l} A\ \text{cubus} \\ =\,B \text{ in } A\ \text{quad. } 2 \end{array}\right\}\ \text{æquabitur } \begin{array}{l} B\ \text{quadrat. in } A\ \text{bis} \\ =\,D\ \text{plano in } A \end{array} \]
\marginal{$\pm$ ; $\pm$}
\note{Les deux signes en « = » sont des doubles traits surchargés ; en regard, dans la marge, deux signes de correction et un long trait qui mène au premier.}

et ductis iisdem in $B$
\[ \left.\begin{array}{l} B \text{ in } A\ \text{quad.} \\ +\,B\ \text{quad. in } A\ \text{bis} \end{array}\right\}\ \text{æquabitur } \begin{array}{l} B\ \text{cubo bis} \\ +\,D\ \text{plano in } B \end{array} \]
Iungantur ducta æqualia æqualibus
\[ \left.\begin{array}{l} A\ \text{cubus} \\ +\,B \text{ in } A\ \text{quadrat. ter} \\ +\,B\ \text{quadrat. in } A\ \text{bis} \end{array}\right\}\ \text{æquabitur } \begin{array}{l} B\ \text{quadrat. in } A\ \text{bis} \\ =\,D\ \text{plano in } A \\ =\,B\ \text{cubo bis} \\ +\,D\ \text{plano in } B \end{array} \]
\marginal{$\pm$ ; $\pm$}
\note{Une tache sur « A cubus ». Deux signes de correction dans la marge, en regard des deux doubles traits du second membre.}

et deleta utrinque adfectione $B$ quadrati in $A$ bis et ad ordinationem æqualitatis translata per antithesin $D$ plani in $A$ adfectione
\[ \left.\begin{array}{l} A\ \text{cubus} \\ =\,B \text{ in } A\ \text{quadrat. ter} \\ -\,D\ \text{plano in } A \end{array}\right\}\ \text{æquabitur } \begin{array}{l} B\ \text{cubo bis} \\ +\,D\ \text{plano in } B \end{array} \]
\marginal{$+$}
\note{Devant « B in A quadrat. ter », un double trait, avec un trait venu de la marge, où un « + » souligné est écrit en regard : c'est le signe de l'énoncé de la proposition II.}

quod quidem ita se habet

$1C + 30Q - 24N$ æquetur $22400$

igitur $1Q + 20N$ æquatur $224$ \struck{et} fit $1N$ $8$.
\marginal{2240}
\note{« 22400 » est souligné, et « 2240 », souligné, est écrit en regard dans la marge de gauche : c'est la correction. Avec $B = 10$ et $D = 24$ : $2000 + 240 = 2240$, $200 + 24 = 224$, et $64 + 160 = 224$ ; $512 + 1920 - 192 = 2240$. Devant « 1C », un petit signe surchargé.}

\subsection*{prop. III}
\note{Après « III », un signe noirci d'encre ; le numéro pourrait avoir été corrigé.}

\[ \text{Si } \left.\begin{array}{l} A\ \text{cubus} \\ =\,B \text{ in } A\ \text{quad. ter} \\ +\,D\ \text{plano in } A \end{array}\right\}\ \text{æquetur } \left\{\begin{array}{l} D\ \text{plano in } B \\ -\,B\ \text{cubo bis} \end{array}\right. \]
\marginal{$-$}
\note{Devant « B in A quad. ter », un signe repris (double trait surchargé), avec un trait venu de la marge, où un « $-$ » est écrit en regard, souligné : c'est le signe de la reprise au bas de la page.}

et sit $B$ quadratum ter maius $D$ plano
\note{« ter » : lecture des lettres ; la condition que demande l'exemple est que $D$ plan passe deux fois $B$ quadratum ($236 > 200$).}
\[ \left.\begin{array}{l} B \text{ in } A\ \text{bis} \\ -\,A\ \text{quadrato} \end{array}\right\}\ \text{æquatur } \begin{array}{l} D\ \text{plano} \\ -\,B\ \text{quadrato } 2 \end{array} \]
\[ \text{quoniam enim } \left.\begin{array}{l} B \text{ in } A\ 2 \\ -\,A\ \text{quadrato} \end{array}\right\}\ \text{æquat. } \begin{array}{l} D\ \text{plano} \\ -\,B\ \text{quadrato } 2 \end{array} \]
ductis omnibus in $B - A$
\[ \left.\begin{array}{l} B\ \text{quad. in } A\ \text{bis} \\ -\,B \text{ in } A\ \text{quad.} \\ -\,B \text{ in } A\ \text{quad. } 2 \\ +\,A\ \text{cubo} \end{array}\right\}\ \text{æquabit. } \begin{array}{l} D\ \text{plano in } B \\ -\,B\ \text{cubo bis} \\ -\,D\ \text{plano in } A \\ +\,B\ \text{quadrat. in } A\ \text{bis} \end{array} \]
et ordinata æqualitate
\[ \left.\begin{array}{l} A\ \text{cubus} \\ -\,B \text{ in } A\ \text{quad. ter} \\ +\,D\ \text{plano in } A \end{array}\right\}\ \text{æquabitur } \begin{array}{l} D\ \text{plano in } B \\ =\,B\ \text{cubo bis} \end{array} \]
\note{Un trait horizontal traverse la marge de gauche au-dessus de « et ordinata æqualitate ». Une tache sur le « D » de « D plano in A ».}

quod quidem ita se habet

$1C - 30Q + 236N$ æquetur $360$

Igitur $20N - 1Q$ æquat. $36$ et fit $1N$ $2$ vel $18$.
\marginal{vel 18}
\note{« vel 18. » est ajouté au bout de la ligne, souligné, et répété, souligné aussi, dans la marge de gauche. L'exemple est juste : avec $B = 10$ et $D = 236$, $2360 - 2000 = 360$, $236 - 200 = 36$, et $A^2 - 20A + 36 = 0$ a pour racines $2$ et $18$. La page s'arrête ici ; le bas du feuillet est blanc.}

\page{132}
\note{Feuillet numéroté « 62 » à l'encre en haut à droite. Au-dessus, la même bande du feuillet numéroté « 63 » qu'à la vue 130, qui n'est pas de cette page. À gauche, le bord du feuillet précédent. La page continue la proposition III du chapitre X (vue 131) : elle en tire la seconde racine $B$.}

Iisdem positis ipsa $A$ fit quoque $B$ sive $B$ triplum quadratum præstet sive cedat $D$ plano

quoniam enim proponitur
\[ \left.\begin{array}{l} A\ \text{cubus} \\ -\,B \text{ in } A\ \text{quad. } 3 \\ +\,D\ \text{plano in } A \end{array}\right\}\ \text{æquatur } \begin{array}{l} D\ \text{plano in } B \\ -\,B\ \text{cubo bis} \end{array} \]
\marginal{æquari}
\note{« æquatur » est souligné, et « æquari » est écrit en regard dans la marge de gauche.}

Ipsa autem $A$ fit quoque $B$
\[ \text{Igitur } \left.\begin{array}{l} B\ \text{cubus} \\ -\,B\ \text{cubo ter} \\ +\,D\ \text{plano in } B \end{array}\right\}\ \text{æquabit. } \begin{array}{l} D\ \text{plano in } B \\ -\,B\ \text{cubo bis} \end{array} \]
quod quidem ita se habet

$1C - 30Q + 236N$ æquat. $360$

ostensa est $1N$ $2$ vel $18$ eadem quoque est $10$

$1C - 30Q + 264N$ æquatur $640$ fit $1N$ $4, 10, 16$.
\marginal{4. vel. 16.}

nam $20N - 1Q$ æquatur $64$ et fit $1N$ $4$ vel $16$.
\note{« 4, 10, 16. » est souligné, les trois nombres écrits presque sans séparation (« 41,016 ») ; la marge de gauche porte en regard « 4. vel. 16. », souligné. Les exemples sont justes : $10$ satisfait $1C - 30Q + 236N = 360$ ($1000 - 3000 + 2360$) ; avec $D = 264$, $2640 - 2000 = 640$, et $A^2 - 20A + 64 = 0$ a pour racines $4$ et $16$, la troisième étant $B = 10$.}

\subsection*{prop. IIII}

\[ \text{Si } \left.\begin{array}{l} B \text{ in } A\ \text{quadratum} \\ +\,D\ \text{plano in } A \\ -\,A\ \text{cubo} \end{array}\right\}\ \text{æquetur } \begin{array}{l} B\ \text{cubo bis} \\ +\,D\ \text{plano in } B \end{array} \]
\marginal{ter}
\note{« quadratum » est souligné, et « ter », souligné, est écrit en regard dans la marge de gauche : $B$ in $A$ quadratum ter, comme dans la reprise au bas de la page.}
\[ \left.\begin{array}{l} A\ \text{quadrat.} \\ -\,B \text{ in } A\ \text{bis} \end{array}\right\}\ \text{æquatur } \begin{array}{l} B\ \text{quadrato bis} \\ +\,D\ \text{plano} \end{array} \]
\[ \text{quoniam enim } \left.\begin{array}{l} A\ \text{quadrat.} \\ -\,B \text{ in } A\ \text{bis} \end{array}\right\}\ \text{æquat. } \begin{array}{l} B\ \text{quadrato bis} \\ +\,D\ \text{plano} \end{array} \]
ductis omnibus in $B - A$
\[ \left.\begin{array}{l} B \text{ in } A\ \text{quadratum} \\ -\,B\ \text{quad. in } A\ \text{bis} \\ -\,A\ \text{cubo} \\ +\,B \text{ in } A\ \text{quad. } 2 \end{array}\right\}\ \text{æquabitur } \begin{array}{l} B\ \text{cubo bis} \\ =\,B \text{ in } D\ \text{plano} \\ =\,B\ \text{quad. in } A\ \text{bis} \\ -\,D\ \text{plano in } A \end{array} \]
\marginal{$\pm$}
\note{Devant « B in D plano » et « B quad. in A bis », des doubles traits surchargés ; le « B » de « B in D plano » est repris. En regard, dans la marge de gauche, un signe de correction (une croix sur un double trait). Le produit demande $+ B$ in $D$ plano et $- B$ quad. in $A$ bis.}

et æqualitate ordinata
\[ \left.\begin{array}{l} B \text{ in } A\ \text{quad. ter} \\ +\,D\ \text{plano in } A \\ -\,A\ \text{cubo} \end{array}\right\}\ \text{æquabit. } \begin{array}{l} B\ \text{cubo bis} \\ +\,D\ \text{plano in } B \end{array} \]
quod quidem ita se habet

$30Q + 24N - 1C$ æquatur $2240$ igitur $1Q - 20N$ æquatur $224$ et fit $1N$ $28$.
\note{L'exemple est juste : avec $B = 10$ et $D = 24$, $2000 + 240 = 2240$, $200 + 24 = 224$, et $784 - 560 = 224$ ; $23520 + 672 - 21952 = 2240$.}

Iisdem \uncertain{expositis} fit $A$ \struck{\ill{}} quoque $B$.
\note{Le deuxième mot commence par une lettre reprise (« e » sur « p » ?) ; le mot barré après « A » est noirci (peut-être « æqualis »). La page s'arrête sur cette phrase, qui annonce la seconde racine ; le reste du feuillet est blanc.}

\page{133}
\note{Verso du feuillet numéroté 62 (vue 132) : aucun numéro en tête. Le feuillet est plus petit que la vue ; à droite, le bord d'autres feuillets, avec des bouts de lignes qui ne sont pas de cette page. La page continue la proposition IIII du chapitre X : elle démontre la seconde racine annoncée à la fin de la vue 132.}

quoniam enim proponitur
\[ \left.\begin{array}{l} B \text{ in } A\ \text{quad. ter} \\ =\,D\ \text{plano in } A \\ -\,A\ \text{cubo} \end{array}\right\}\ \text{æquari } \begin{array}{l} B\ \text{cubo bis} \\ +\,D\ \text{plano in } B \end{array} \]
\marginal{$\pm$}
\note{Devant « D plano in A », un double trait surchargé, avec un trait venu de la marge, où un signe de correction (une croix sur un double trait) est écrit en regard, souligné ; l'énoncé de la proposition IIII porte $+ D$ plano in $A$. En haut à droite, une petite croix.}

ipsa autem $A$ fit quoque $B$
\[ \text{Igitur } \left.\begin{array}{l} B\ \text{cubus ter} \\ +\,D\ \text{plano in } B \\ -\,B\ \text{cubo} \end{array}\right\}\ \text{æquabit. } \begin{array}{l} B\ \text{cubo bis} \\ +\,D\ \text{plano in } B \end{array} \]
quod quidem ita se habet

$30Q + 24N - 1C$ æquatur $2240$ fit $1N$ $10$.
\note{L'exemple est juste : $3000 + 240 - 1000 = 2240$.}

\subsection*{prop. V.}

\[ \text{Si } \left.\begin{array}{l} B \text{ in } A\ \text{quad. } 3 \\ -\,D\ \text{plano in } A \\ -\,A\ \text{cubo} \end{array}\right\}\ \text{æquabitur } \begin{array}{l} B\ \text{cubo bis} \\ -\,D\ \text{plano in } B \end{array} \]
\marginal{æquetur}
\note{« æquabitur » est souligné, et « æquetur », souligné, est écrit en regard dans la marge de gauche. Un petit trait oblique à droite de la ligne.}
\[ \left.\begin{array}{l} A\ \text{quad.} \\ -\,B \text{ in } A\ \text{bis} \end{array}\right\}\ \text{æquabitur } \left\{\begin{array}{l} B\ \text{quad. bis} \\ -\,D\ \text{plano} \end{array}\right. \]
\[ \text{quoniam enim } \left.\begin{array}{l} A\ \text{quadrat.} \\ -\,B \text{ in } A\ \text{bis} \end{array}\right\}\ \text{æquatur } \begin{array}{l} B\ \text{quadrato bis} \\ -\,D\ \text{plano} \end{array} \]
ductis omnibus in $B - A$
\[ \left.\begin{array}{l} B \text{ in } A\ \text{quadratum} \\ -\,B\ \text{quad. in } A\ \text{bis} \\ -\,A\ \text{cubo} \\ +\,B \text{ in } A\ \text{quadrat. bis} \end{array}\right\}\ \text{æquabit. } \begin{array}{l} B\ \text{cubo bis} \\ -\,D\ \text{plano in } B \\ -\,B\ \text{quadrato in } A\ \text{bis} \\ +\,D\ \text{plano in } A \end{array} \]
Et æqualitate ordinata
\[ \left.\begin{array}{l} B \text{ in } A\ \text{quadrat. ter} \\ -\,D\ \text{plano in } A \\ -\,A\ \text{cubo} \end{array}\right\}\ \text{æquabit. } \begin{array}{l} B\ \text{cubo bis} \\ -\,D\ \text{plano in } B. \end{array} \]
quod quidem ita se habet

$30Q - 156N - 1C$ æquatur $440$

Igitur $1Q - 20N$ æquatur $44$ et fit $1N$ $22$
\note{L'exemple est juste : avec $B = 10$ et $D = 156$, $2000 - 1560 = 440$, $200 - 156 = 44$, et $484 - 440 = 44$.}

Iisdem positis fit $A$ quoque $B$

quoniam enim proponitur
\[ \left.\begin{array}{l} B \text{ in } A\ \text{quad. } 3 \\ -\,D\ \text{plano in } A \\ -\,A\ \text{cubo} \end{array}\right\}\ \text{æquari } \begin{array}{l} B\ \text{cubo bis} \\ -\,D\ \text{plano in } B \end{array} \]
Ipsa autem $A$ fit quoque $B$
\[ \text{Igitur } \left.\begin{array}{l} B\ \text{cubus } 3 \\ -\,D\ \text{plano in } B \\ -\,B\ \text{cubo} \end{array}\right\}\ \text{æquabitur } \begin{array}{l} B\ \text{cubo bis} \\ -\,D\ \text{plano in } B. \end{array} \]
\note{Un trait épais et une tache passent sur « autem A fit » ; on ne peut dire si les mots sont biffés. La page s'arrête ici ; la vérification (« quod quidem ita se habet », « $30Q - 156N - 1C$ æquatur $440$ fit $1N$ $10$ ») est en tête du feuillet numéroté 63, dont le haut paraît aux vues 130 et 132 — ce qui vérifie : $3000 - 1560 - 1000 = 440$.}

\page{134}
\note{Feuillet numéroté « 63 » à l'encre en haut à droite, dont le haut paraît déjà aux vues 130 et 132. Il est plus court que le fond : sous son bord inférieur paraissent trois lignes d'un autre feuillet (« \ldots{} continuæ proportionales / Z \struck{summa} differentia trium primarum \struck{a tribus postremis} æquatur / reliquis, a tribus postremis, et fit A secunda. »), qui ne sont pas de cette page ; elles reparaissent à la vue 138, où elles sont transcrites. Les deux premières lignes achèvent la proposition V du chapitre X (vue 133).}

quod quidem ita se habet

$30Q - 156N - 1C$ æquatur $440$ fit $1N$ $10$
\note{« 440 » : le dernier chiffre est écrit comme un 6 ; la vue 133 donne 440, et $3000 - 1560 - 1000 = 440$.}

\subsection*{Prop. VI}

\marginal{A}
Ad quadrato-quadraticam pertinens
\note{Un « A » est écrit devant la ligne, en retrait, comme un signe de renvoi.}

\[ \text{Si } \left.\begin{array}{l} A\ \text{quadrato-quadrat.} \\ -\,X \text{ in } A\ \text{cubum } 2 \\ +\,X\ \text{cubo in } A\ 4 \end{array}\right\}\ \text{æquetur } X\ \text{quadrato-quadrato } 2 \]
\[ \left.\begin{array}{l} X\ \text{quad. in } A\ \text{quad.} \\ -\,A\ \text{quad.-quad.} \end{array}\right\}\ \text{æquatur } \begin{array}{l} X\ \text{quad.-quad.} \\ -\,X\ \text{cub. } 2 \text{ in } L\,X\ \text{quadrati} \end{array} \]
\marginal{2 ; 4 ; æquetur ; 3}
\note{Dans la seconde égalité, « A quad. », « æquatur », « X quad.-quad. » et « X quadrati » sont soulignés ou suivis d'un trait, et la marge de gauche porte en regard, soulignés, « 2 », « 4 » (encadré), « æquetur » et « 3 » : les corrections donnent $X$ quad. in $A$ quad. 2 $-$ $A$ quad.-quad. æquetur $X$ quad.-quad. 4 $-$ $X$ cub. 2 in $L\,X$ quadrati 3, comme dans la note marginale plus bas.}

\[ \text{ergo } \left.\dfrac{\begin{array}{l} X\ \text{quad.-quad. } 4 \\ +\,A\ \text{quad.-quad.} \\ -\,X\ \text{quad. in } A\ \text{quad. } 2 \end{array}}{X\ \text{cub. } 2}\right\}\ \text{æquabitur } L\,X\ \text{quadrati} \]
\marginal{3}
\note{« quadrati » est souligné, et « 3 », souligné, est écrit en regard dans la marge.}

et omnibus quadratis et per $X$ cubo-cubum $4$ ductis adhibitaque congruenti metathesi
\[ \left.\begin{array}{l} X\ \text{cubo-cubus in } A\ \text{quad. } 16 \\ =\,X\ \text{quad. in } A\ \text{cubo-cubum } 4 \\ -\,A\ \text{quad.-quad.} \\ -\,X\ \text{quad.-quad. in } A\ \text{quad.-quad. } 12 \end{array}\right\}\ \text{æquabit. } X\ \dfrac{\text{quad.-quad.-quad.-quad.}}{4} \]
\marginal{$\pm$ ; quad.-quad.}
\note{Devant « X quad. in A cubo-cubum 4 », un double trait surchargé, avec un trait venu de la marge, où un signe de correction (une croix sur un double trait) est écrit en regard ; « A quad.-quad. » est souligné, et « quad. quad. » écrit en regard dans la marge : $A$ quad.-quad.-quad.-quad. Le « 4 » sous le second membre est posé sous le trait, comme un diviseur ou un coefficient. La ligne est laissée telle ; plusieurs de ses coefficients ne répondent pas à l'égalité qui précède.}

quod quidem ita se habet. \ill{} \uncertain{fit} \uncertain{rursus} \uncertain{primam} æquationum
\[ \left.\begin{array}{l} X\ \text{cubus in } A\ 4 \\ -\,X \text{ in } A\ \text{cubum } 2 \end{array}\right\}\ \text{æquatur } \begin{array}{l} X\ \text{quad.-quad. } 2 \\ -\,A\ \text{quad.-quad.} \end{array} \]
tunc æqualitatis parte utraque quadrata omnibusque rite ordinatis in eam ipsam æqualitatem inciditur quadrato-cubicam quæ exposita est
\note{Les quatre mots après « se habet. » sont d'une écriture courante que les coupes n'ont pas résolue ; seul « æquationum » est sûr. La première égalité de la proposition, ses termes répartis de part et d'autre, élevée au carré, redonne l'équation du sixième degré qui précède (« quadrato-cubicam »).}

Itaque $A$ quad. fit $X$ quadrat. plus minusve latere binomii residui seu negatæ $L\,X$ quad.-quad. $12 - X$ quad. $3$

fit $X$. $1$ $A$ $1N$.
\[ \left.\begin{array}{l} 1QQ \\ -\,2C \\ +\,4N \end{array}\right\}\ \text{æquatur } 2 \]
Igitur $2Q - 1QQ$ æquabitur $\overline{4 - L\,12}$

et $1Q$ fit $1$ plus minusve latere binomii negatæ $L\,12 - 3$.
\note{Un trait est tiré au-dessus de « $4 - L\,12$ ». L'exemple suit la règle corrigée : avec $X = 1$, $2A^2 - A^4 = 4 - \sqrt{12}$, d'où $A^2 = 1 \pm \sqrt{\sqrt{12} - 3}$. Dans la marge de gauche, en regard, un signe de correction fait de trois traits.}

\marginal{quoniam enim $\left.\begin{array}{l} Xq \text{ in } Aq.\ 2 \\ -\,Aqq \end{array}\right|\begin{array}{l} Xqq.\ 4 \\ -\,Xc\ 2 \text{ in } L\,Xq\,3. \end{array}$ ergo $Aq. \,|\, Xq.$ \struck{\ill{}} minus plusve $L \left\{\begin{array}{l} Xqq - Xqq\ 4 \\ +\,Xc\ 2 \text{ in } L\,Xq\,3. \end{array}\right.$ id est ex lege Logisticæ $L.\,Xqq.\,12 - Xq\,3$.}
\note{Note marginale d'une encre plus noire, dans la marge de gauche, en regard de la règle et de l'exemple : elle refait le calcul sur l'égalité corrigée. Le trait vertical y tient lieu de « æquatur ». Le mot barré après « Xq. » est noirci. Le calcul est juste pour les nombres : $X^4 - 4X^4 + 2\sqrt{3}\,X^4 = (\sqrt{12} - 3)X^4$ ; l'écriture abrégée « L. Xqq. 12 $-$ Xq 3 » ne garde pas l'homogénéité, et elle est laissée telle.}

\page{135}
\note{Verso du feuillet numéroté 63 (vue 134) : aucun numéro en tête. Le feuillet est plus court que le fond : sous son bord inférieur paraît la fin de la vue 133 (« Ipsa autem A fit quoque B / Igitur B cubus 3 \ldots{} æquabitur B cubo bis $-$ D plano in B »), qui n'est pas de cette page. À droite, le bord d'un autre feuillet. La page ouvre un nouveau chapitre ; la proposition VI de la vue 134 s'y achève sans suite.}

\subsection*{Singularium aliquot constitutionum ad æqualitates multipliciter affectas pertinentium Collectio. Cap. \struck{X}XI}
\marginal{XI}
\note{Le numéro du chapitre est d'abord écrit « XXI », le premier « X » barré, et souligné ; « XI », souligné, est répété dans la marge de gauche. « affectas » est souligné.}

\subsection*{propositio 1}

\[ \text{Si } \left.\begin{array}{l} A\ \text{quadrat.} \\ +\,B \text{ in } A \end{array}\right\}\ \text{æquat. } B\ \text{quadrato} \]
est $B$ dupla longitudo secta in tria proportionalia segmenta quorum primum idemque maius est $B$. secundum $A$. quo in opere dicitur $B$ secari media et extrema ratione \struck{simul} simpla.
\marginal{fit prima proportionalis ipsa longitudo simpla $B$ secunda $A$ tertia est.}
\marginal{semel}
\note{« est B dupla » : « dupla » se rapporte à « longitudo », la longueur double de $B$. « simul » est souligné et barré, et « semel » est écrit en regard dans la marge ; le mot qui suit se lit « simpla » ou « sumpta ». La note marginale du haut, de la même main, en trois lignes, reprend la définition. Les trois segments sont $B$, $A$ et $B - A$, et $B : A = A : (B - A)$ donne bien $A^2 + BA = B^2$.}

\subsection*{prop. II}

\[ \text{Si } \left.\begin{array}{l} A\ \text{cubus} \\ +\,B \text{ in } A\ \text{quadrat.} \\ +\,B\ \text{quadrato in } A \end{array}\right\}\ \text{æquetur } B\ \text{cubo} \]
est $B$ dupla longitudo secta in quatuor continue proportionalia segmenta quorum primum idemque maius est $B$, secundum $A$. quo in opere dicitur $B$ secari media et extrema ratione duplicata

\subsection*{prop. III}

\[ \text{Si } \left.\begin{array}{l} A\ \text{quad.-quad.} \\ +\,B \text{ in } A\ \text{cubum} \\ +\,B\ \text{quad. in } A\ \text{quad.} \\ +\,B\ \text{cubo in } A \end{array}\right\}\ \text{æquetur } B\ \text{quad.-quad.} \]
est $B$ dupla longitudo secta in quinque \add{continue} proportionalia segmenta quorum primum idemque maius est $B$ secundum $A$ quo in opere dicitur $B$ secari media et extrema ratione triplicata
\marginal{continue}
\note{Entre « quinque » et « proportionalia », un signe d'insertion ; « continue », souligné, est écrit en regard dans la marge de gauche.}

\subsection*{prop. IIII}

\[ \text{Si } \left.\begin{array}{l} A\ \text{quadrato-cubus} \\ +\,B \text{ in } A\ \text{quad.-quad.} \\ +\,B\ \text{quad. in } A\ \text{cubum} \\ +\,B\ \text{cubo in } A\ \text{quad.} \\ +\,B\ \text{quad.-quad. in } A \end{array}\right\}\ \text{æquetur } B\ \text{quadrato-cubo} \]
est $B$ dupla longitudo secta in sex proportionalia segmenta quorum primum idemque maius est $B$ secundum $A$. quo in opere dicitur $B$ secari media et extrema ratione quadruplicata

\subsection*{prop. V}

\[ \text{Si } \left.\begin{array}{l} A\ \text{cubo-cubus} \\ +\,B \text{ in } A\ \text{quadrato-cubum} \\ +\,B \text{ in } A\ \text{quad.-quad.} \\ +\,B\ \text{cubo in } A\ \text{cubum} \\ +\,B\ \text{quad.-quad. in } A\ \text{quad.} \\ +\,B\ \text{quadrato-cubo in } A \end{array}\right\}\ \text{æquetur } B\ \text{cubo-cubo} \]
\note{À la deuxième ligne, « B in A quadrato-cubum » ; la troisième porte « B in A quad.-quad. », là où la suite des propositions précédentes demande $B$ quad. in $A$ quad.-quad. La page s'arrête sur cet énoncé, au bas du feuillet, sans la conclusion ; la conclusion est en tête du feuillet suivant, vue 136.}

\page{136}
\note{Feuillet numéroté « 64 » à l'encre en haut à droite. Il est plus court que le fond : sous son bord inférieur paraissent, sur le feuillet de dessous, les mêmes lignes qu'à la vue 134 (« \ldots{} continuæ proportionales / Z \struck{summa} differentia trium primarum \ldots{} », suivies à droite de nombres coupés par le bord, « \ldots{} 4. 10. 10. 52. »), qui ne sont pas de cette page ; les lignes sont transcrites à la vue 138. À gauche, le bord du feuillet précédent, avec des fins de lignes. Les trois premières lignes achèvent la proposition V du chapitre XI (vue 135).}

est $B$ dupla longitudo secta in septem continue proportionalia segmenta quorum primum idemque maius est $B$ secundum $A$ quo in opere dicitur $B$ secari media et extrema ratione quintuplicata.

\subsection*{Earundem collectio altera Cap. \struck{X}XII}
\marginal{XII}
\note{Le premier « X » du numéro est noyé sous une tache ; « XII », souligné, est répété dans la marge de gauche.}

\subsection*{propositio 1}

\[ \text{Si } \left.\begin{array}{l} A\ \text{quadrat.} \\ +\,B \text{ in } A \end{array}\right\}\ \text{æquetur } B \text{ in } Z \]
est $B$ prima minor inter extremas in serie trium proportionalium aggregatum vero reliquarum duarum est $Z$ et fit $A$ secunda
\note{L'énoncé est juste : si $B$, $A$, $C$ sont proportionnelles et $A + C = Z$, alors $A^2 = B(Z - A)$.}

\subsection*{prop. II}

\[ \text{Si } \left.\begin{array}{l} A\ \text{cubus} \\ +\,B \text{ in } A\ \text{quadratum} \\ +\,B\ \text{quadrato in } A \end{array}\right\}\ \text{æquetur } B\ \text{quadrat. in } Z \]
est $B$ prima in serie quatuor continue proportionalium aggregatum vero reliquarum trium $Z$ et fit $A$ secunda
\note{Une tache couvre le milieu de « quadratum » dans l'énoncé.}

\subsection*{prop. III}

\[ \text{Si } \left.\begin{array}{l} A\ \text{quadrato-quadrat.} \\ +\,B \text{ in } A\ \text{cubum} \\ +\,B\ \text{quad. in } A\ \text{quad.} \\ +\,B\ \text{cubo in } A \end{array}\right\}\ \text{æquetur } B\ \text{cubo in } Z \]
est $B$ prima in serie quinque continue proportionalium aggregatum vero reliquarum quatuor est $Z$ et fit $A$ secunda

\subsection*{prop. IIII}

\[ \text{Si } \left.\begin{array}{l} A\ \text{quadrato-cubus} \\ +\,B \text{ in } A\ \text{quad.-quad.} \\ +\,B\ \text{quad. in } A\ \text{cub.} \\ +\,B\ \text{cubo in } A\ \text{quad.} \\ +\,B\ \text{quad.-quad. in } A \end{array}\right\}\ \text{æquetur } B\ \text{quad.-quad. in } Z \]
est $B$ prima in serie sex continue proportionalium aggregatum vero reliquarum quinque est $Z$ et fit $A$ secunda.
\note{La page s'arrête ici, au milieu du feuillet ; le reste est blanc.}

\page{137}
\note{Verso du feuillet numéroté 64 (vue 136) : aucun numéro en tête. Le feuillet est plus court que le fond : sous son bord inférieur reparaît la fin de la vue 133 (« Ipsa autem A fit quoque B / Igitur B cubus 3 \ldots »), qui n'est pas de cette page. À droite, un onglet. La page poursuit le chapitre XII de la vue 136.}

\subsection*{prop. V}

\[ \text{Si } \left.\begin{array}{l} A\ \text{cubo-cubus} \\ +\,B \text{ in } A\ \text{quadrato-cubum} \\ +\,B\ \text{quad. in } A\ \text{quad.-quad.} \\ +\,B\ \text{cubo in } A\ \text{cubum} \\ +\,B\ \text{quad.-quad. in } A\ \text{quad.} \\ +\,B\ \text{quadrato-cubo in } A \end{array}\right\}\ \text{æquetur } B\ \text{quadrato-cubo in } Z \]
est $B$ prima in serie septem continue proportionalium aggregatum vero reliquarum sex est $Z$ et fit $A$ secunda
\note{Le « A » de « in A quad. » (quatrième terme) est noirci.}

\subsection*{Earundem collectio tertia Cap. \struck{X}XIII}
\marginal{XIII ; prop. 1.}
\note{« XXIII » est souligné, le premier « X » peut-être biffé ; « XIII », souligné, et « prop. 1. » sont écrits en regard dans la marge de gauche.}

\[ \text{Si } \left.\begin{array}{l} B \text{ in } A \\ -\,A\ \text{quadrato} \end{array}\right\}\ \text{æquetur } B \text{ in } Z \]
est $B$ prima maior inter extremas in serie trium proportionalium differentia vero duarum reliquarum est $Z$ et fit $A$ secunda. potest autem esse duplex. nam est etiam differentia inter primam et secundam.

\[ \text{Sed si } \left.\begin{array}{l} A\ \text{quadratum} \\ -\,B \text{ in } A \end{array}\right\}\ \text{æquetur } B \text{ in } Z \]
est $B$ prima minor inter extremas differentia vero duarum reliquarum est $Z$ et fit $A$ similiter secunda

Sunto proportionales $4.\ 6.\ 9$

$9N - 1Q$ æquatur $18$ fit $1N$ $6$ atque etiam $3$

Sed si unum quadratum $- 4N$ æquetur $12$ fit $1N$ $6$.
\note{Une tache sur la seconde lettre de « maior inter extremas ». « unum » est écrit « vnu » avec un trait. Les exemples sont justes : avec $B = 9$ et $Z = 6 - 4 = 2$, $9A - A^2 = 18$ a pour racines $6$ et $3$ ($3 = 9 - 6$, la « différence entre la première et la seconde ») ; avec $B = 4$ et $Z = 9 - 6 = 3$, $A^2 - 4A = 12$ donne $6$.}

\subsection*{prop. II}

\[ \text{Si } \left.\begin{array}{l} A\ \text{cubus} \\ =\,B \text{ in } A\ \text{quad.} \\ +\,B\ \text{quad. in } A \end{array}\right\}\ \text{æquetur } B\ \text{quadrato in } Z \]
\marginal{$-$}
\note{Le signe devant « B in A quad. » est un double trait surchargé, avec un long trait venu de la marge de gauche, où un « $-$ » est écrit en regard : c'est le signe que demandent la règle et les exemples. « Si A cubus » est souligné.}
est $B$ prima in serie quatuor continue proportionalium differentia vero trium reliquarum alterne sumptarum est $Z$ et fit $A$ secunda

Sunto proportionales continue $\overset{\text{I}}{1}.\ \overset{\text{II}}{2}.\ \overset{\text{III}}{4}.\ \overset{\text{IIII}}{8}.$

$1C - 8Q + 64N$ æquatur $192$ fit $1N$ $4$.

vel $1C - 1Q + 1N$ æquatur $6$ fit $1N$ $2$
\note{Au-dessus des quatre proportionnelles, des traits verticaux en nombre croissant les numérotent. Un trait est tiré au-dessus du premier exemple, qui est souligné en partie. Les deux exemples sont justes : la série lue à rebours, $B = 8$, $A = 4$, $Z = 4 - 2 + 1 = 3$, donne $64 - 128 + 256 = 192 = 64 \cdot 3$ ; la série lue dans l'ordre, $B = 1$, $A = 2$, $Z = 2 - 4 + 8 = 6$, donne $8 - 4 + 2 = 6$.}

\marginal{Si $\left.\begin{array}{l} Bc \text{ in } A \\ -\,Bq \text{ in } Aq \\ +\,B \text{ in } Ac \\ -\,Aqq \end{array}\right| Bc \text{ in } Z.$ est $B$ prima e serie quinque proportionalium $Z$ differentia quatuor reliquarum alterne sumptarum et fit $A$ differentia primæ et secundæ.}
\note{Note marginale d'une encre plus noire, en bas de la marge de gauche, sa dernière ligne courant sous le texte sur toute la largeur de la page : l'énoncé suivant de la série, pour cinq proportionnelles. Chacun des trois premiers termes est suivi d'un petit signe, lu « q », « 3 », « 3 », que le transcripteur ne sait rattacher ni au terme ni à un coefficient ; la lecture des puissances est donnée d'après le patron des propositions I et II, et elle est douteuse. Le trait vertical tient lieu de « æquetur ». La fin, « differentia primæ et secundæ », est nette, alors que les propositions I et II disent « fit A secunda ».}

\page{138}
\note{Feuillet numéroté « 65 » à l'encre en haut à droite. Au-dessus, sur une bande, le haut d'un autre feuillet, numéroté « 67 », avec un titre qui n'est pas de cette page : « AD Tabulas Prostaphæreseôn Lunarium / Propositio 1 / Invenire primam Prostaphæresim » ; voir la vue 140. Le feuillet 65 poursuit le chapitre XIII (vue 137). Ses marges sont couvertes de notes d'une encre plus noire, de la main des corrections, qui énoncent d'autres propositions de la même série ; c'est leur fin que l'on voyait sous le bord des feuillets 63 et 64 aux vues 134 et 136.}

\subsection*{prop. III}

\[ \text{Si } \left.\begin{array}{l} B\ \text{cubus in } A \\ -\,B\ \text{quad. in } A\ \text{quad.} \\ +\,B \text{ in } A\ \text{cubum} \\ -\,A\ \text{quad.-quad.} \end{array}\right\}\ \text{æquetur } B\ \text{cubo in } Z \]
est $B$ prima maior in serie quinque continue proportionalium differentia vero quatuor reliquarum sumptarum alterne est $Z$ et fit $A$ secunda maior inter medias, potest autem esse duplex --- X°
\note{« X° » : un signe de renvoi, fait d'un X surmonté d'un petit o, qui revient deux fois dans la marge de gauche et une fois en haut, près de l'encadré de droite.}

\[ \text{Sed si } \left.\begin{array}{l} A\ \text{quad.-quad.} \\ -\,B \text{ in } A\ \text{cubum} \\ +\,B\ \text{quad. in } A\ \text{quadrat.} \\ -\,B\ \text{cubo in } A \end{array}\right\}\ \text{æquetur } B\ \text{cubo in } Z \]
est $B$ prima minor inter extremas differentia vero quatuor reliquarum sumptarum alterne est $Z$ et fit $A$ secunda minor \struck{inter} \add{inter} medias
\note{« inter » est barré et récrit au-dessus de la ligne.}

Sunto proportionales continue $\overset{\text{I}}{1}.\ \overset{\text{II}}{2}.\ \overset{\text{III}}{4}.\ \overset{\text{IIII}}{8}.\ \overset{\text{V}}{16}.$
\[ \left.\begin{array}{l} 4096\,N \\ -\,256\,Q \\ +\,16\,C \\ -\,1\,QQ \end{array}\right\}\ \text{æquatur } 20{,}480 \text{ fit } 1N\ 8. \]
\[ \text{vel } \left.\begin{array}{l} 1\,QQ \\ -\,1\,C \\ +\,1\,Q \\ -\,1\,N \end{array}\right\}\ \text{æquatur } 10 \text{ fit } 1N\ 2. \]
\note{Les exemples sont justes : la série lue à rebours, $B = 16$, $Z = 8 - 4 + 2 - 1 = 5$, donne $32768 - 16384 + 8192 - 4096 = 20480 = 4096 \cdot 5$ ; la série lue dans l'ordre, $B = 1$, $Z = 16 - 8 + 4 - 2 = 10$, donne $16 - 8 + 4 - 2 = 10$. Le « 2 » à droite de « 256 » est le « Q » mal formé ou un reste de correction.}

\subsection*{prop. IIII}

\[ \text{Si } \left.\begin{array}{l} A\ \text{quadrato-cubus} \\ -\,B \text{ in } A\ \text{quad.-quad.} \\ +\,B\ \text{quad. in } A\ \text{cubum} \\ -\,B\ \text{cubo in } A\ \text{quad.} \\ +\,B\ \text{quad.-quad. in } A \end{array}\right\}\ \text{æquetur } B\ \text{quad.-quad. in } Z. \]
est $B$ prima \add{maior} in serie sex continue proportionalium differentia vero quinque reliquarum alterne sumptarum est $Z$ et fit $A$ secunda

Sunto proportionales continue $\overset{\text{I}}{1}.\ \overset{\text{II}}{2}.\ \overset{\text{III}}{4}.\ \overset{\text{IIII}}{8}.\ \overset{\text{V}}{16}.\ \overset{\text{VI}}{32}.$
\marginal{maior}
\note{« maior » est écrit au-dessus de « prima », et répété, souligné, dans la marge de gauche. La page s'arrête sur la série, sans exemple chiffré ; le bas du feuillet porte seulement la fin de la dernière note marginale.}

\marginal{Si $\left.\begin{array}{l} Bc \text{ in } A \\ +\,Bq \text{ in } Aq\ 2 \\ -\,B \text{ in } Z \text{ in } Aq\ 2 \\ +\,B \text{ in } Ac \\ -\,Aqq \end{array}\right\}$ æqu. $Bq \text{ in } Zq$. fit $B$ prima maior in eadem serie $A$ vero tertia $Z$ autem differentia quatuor primarum alterne sumptarum. X°}
\note{Première note de la marge de gauche, en regard de la proposition III.}

\marginal{\struck{\ill{} \ill{} \ill{} / \ill{} \ill{} \ill{} / \ill{} \ill{} \ill{}} L. plano-plani facti ex aggregato duarum reliquarum alterne \add{ambas} sumptarum in maiorem inter medias ut $2$ et $8$. $10$ in $8$ fit $80$. $1$ in $4$. et $4$ in $80$ fit $320$ cuius $L$ est quæsita}
\note{Deuxième note, sous un trait : trois lignes biffées de traits croisés et illisibles, puis cinq lignes qui ne le sont pas ; « ambas » est écrit au-dessus de « sumptarum ». Elle semble donner, sur la série 1, 2, 4, 8, 16, la valeur de la troisième proportionnelle de la note précédente ; $\sqrt{320}$ n'est pourtant pas un terme de la série, et le calcul n'est pas vérifié ici.}

\marginal{$\left.\begin{array}{l} Bcc \text{ in } Aq \\ +\,Bq \text{ in } Acc \\ -\,Bqq \text{ in } Aqq \\ -\,Acc \end{array}\right\}$ æq. $Bqq \text{ in } Z$ fit $B$ prima \add{maior} e serie 5 continue proportionalium $Z$ differentia reliquarum alterne sumptarum et $A$ latus plani facti sub prima maiore et secunda ei proxima.}
\note{Troisième note, sous un trait, en regard des exemples de la proposition III. Les puissances « Bcc », « Acc » sont lues d'après la forme des lettres, qui se lisent aussi « Brr », « Arr » ; la lecture est douteuse et la note n'est pas vérifiée.}

\marginal{Si $\left.\begin{array}{l} Bc \text{ in } A \\ +\,Bq \text{ in } Aq \\ +\,Aqq \end{array}\right| Bc \text{ in } Z$ fit $B$ prima maior in serie quinque continue proportionalium $Z$ summa quatuor posteriorum est $A$ secunda}
\note{Quatrième note, en regard de la proposition IIII. Il y manque le terme $+ B$ in $Ac$ que demande la somme des quatre derniers termes.}

\marginal{Si $\left.\begin{array}{l} Bc \text{ in } A \\ -\,Aqq \end{array}\right\} Bc \text{ in } Z$ est $B$ prima maior et ipsa in serie quinque continue proportionalium $Z$ \struck{summa} differentia trium primarum \struck{a tribus postremis} æquatur reliquis, a tribus postremis, et fit $A$ secunda.}
\note{Cinquième note, sous un trait, dont les deux dernières lignes courent sous le texte sur toute la largeur du pied de la page ; ce sont elles qu'on voyait sous les feuillets 63 et 64 aux vues 134 et 136.}

\marginal{\uncertain{Æquipergata} ad \ill{} sub $A$ in $L$. æquatio sic se habebit ex hypothesi $\left.\begin{array}{l} Bc \text{ in } E\ 40 \\ -\,Bq \text{ in } Eq\ 10 \\ -\,Eqq \end{array}\right|\begin{array}{l} Bc \text{ in } Z \\ -\,Bqq\ 51 \end{array}$ unde ex opere \ill{} \uncertain{parapleroseos} dabitur latus \ill{}}
\note{En haut à droite, en regard du titre de la proposition III, une note de la même encre noire, dans un encadré tracé à la plume qui se referme sur « Bqq 51 » ; ses mots sont serrés contre le bord et plusieurs ne se lisent pas. La lecture des coefficients (« 40 », « 10 », « 51 ») est sûre ; celle des mots est douteuse, et le sens n'est pas établi.}

\page{139}
\note{Verso du feuillet numéroté 65 (vue 138) : aucun numéro en tête ; en haut à gauche, « 65 » vu en miroir, transparence du recto. À droite, le bord d'un autre feuillet. La page s'ouvre sur les exemples de la proposition IIII du chapitre XIII, dont la série (1, 2, 4, 8, 16, 32) clôt la vue 138.}

\[ \left.\begin{array}{l} 1\,QC \\ -\,32\,QQ \\ +\,1024\,C \\ -\,32768\,Q \\ +\,1048576\,N \end{array}\right\}\ \text{æquatur } 11{,}534{,}336 \text{ fit } 1N\ 16. \]
vel $1QC - 1QQ + 1C - 1Q + 1N$ æquatur $22$ et fit $1N$ $2$.
\note{« 1048576 » est écrit « 1048, 576 ». Les exemples sont justes : avec $B = 32$ et $Z = 16 - 8 + 4 - 2 + 1 = 11$, $1048576 \cdot 11 = 11534336$, et $16$ satisfait l'équation ; avec $B = 1$ et $Z = 2 - 4 + 8 - 16 + 32 = 22$, $32 - 16 + 8 - 4 + 2 = 22$.}

Et hæc singula suam habent ex \uncertain{Zetesi} demonstrationem. \struck{Atque} \add{quæ} iam sequitur collectio sua analytico examini subjicienda libere relinquit theoremata. pertinent autem ad æqualitates de multiplicibus radicibus mire explicabiles.
\note{« Atque » paraît biffé, et « quæ » est écrit au-dessus ; les deux premiers mots de l'alinéa sont traversés d'un trait, qui semble être celui du « t » de « et » et n'est pas rendu comme une rature. « Zetesi » est lu d'après les lettres (Z, e, t, e, s, i).}

\subsection*{Collectio quarta. Cap. XXIIII}
\marginal{XIIII}
\note{« XXIIII » est souligné, et « XIIII », souligné, est écrit en regard dans la marge de gauche, comme aux chapitres précédents.}

\subsection*{prop. 1}

\[ \text{Si } \left.\begin{array}{l} B + D \text{ in } A \\ -\,A\ \text{quadrato} \end{array}\right\}\ \text{æquetur } B \text{ in } D. \]
$A$ explicabilis est de qualibet illarum duarum $B$ vel $D$.
\[ \left.\begin{array}{l} 3\,N \\ -\,1\,Q \end{array}\right\}\ \text{æquetur } 2 \text{ fit } 1N\ 1.\ \text{vel } 2. \]
\note{Le « ex » de « explicabilis » est surchargé. L'exemple est juste : $3 - 1 = 2$ et $6 - 4 = 2$.}

\subsection*{prop. II}

\[ \text{Si } \left.\begin{array}{l} A\ \text{cubus} \\ -\,\left.\begin{array}{l} B \\ D \\ G \end{array}\right\} \text{in } A\ \text{quadratum} \\ +\,\left.\begin{array}{l} B \text{ in } D \\ B \text{ in } G \\ D \text{ in } G \end{array}\right\} \text{in } A \end{array}\right\}\ \text{æquetur } B \text{ in } D \text{ in } G. \]
$A$ explicabilis est de qualibet illarum trium $B$, $D$, vel $G$.

$1C - 6Q + 11N$ æquatur $6$

fit $1N$ \struck{\ill{}} $1.\ 2.$ vel $3$.

\struck{prop. III}
\note{Devant « 1. », un chiffre surchargé et biffé. L'exemple est juste : $(x - 1)(x - 2)(x - 3) = x^3 - 6x^2 + 11x - 6$. Le titre « prop. III » est biffé, et le reste de la page est blanc : la proposition n'a pas été écrite ici.}

\page{140}
\note{Feuillet numéroté « 66. » à l'encre en haut à droite. Au-dessus, sur une bande, le haut du feuillet numéroté « 67 », déjà vu à la vue 138, avec son titre (« AD Tabulas Prostaphæreseôn Lunarium / Propositio 1 / Invenire primam Prostaphæresim »), qui n'est pas de cette page. À gauche, le bord du feuillet précédent. La page poursuit le chapitre XIIII de la vue 139.}

\subsection*{prop. III}

\[ \text{Si } \left.\begin{array}{l} \left.\begin{array}{l} B \text{ in } D \text{ in } G \\ B \text{ in } D \text{ in } H \\ B \text{ in } G \text{ in } H \\ D \text{ in } G \text{ in } H \end{array}\right\} \text{in } A \\[1ex] -\left\{\begin{array}{l} B \text{ in } D \\ B \text{ in } G \\ B \text{ in } H \\ D \text{ in } G \\ D \text{ in } H \\ G \text{ in } H \end{array}\right\} \text{in } A\ \text{quadrat.} \\[1ex] +\left.\begin{array}{l} B \\ D \\ G \\ H \end{array}\right\} \text{in } A\ \text{cubum} \end{array}\right\}\ \text{æquetur } B \text{ in } D \text{ in } G \text{ in } H \]
\marginal{$-\,Aqq.$}
\note{Sous le dernier groupe, un trait ; la marge de gauche porte en regard « $-$ Aqq. », le terme $- A$ quad.-quad. qui manque à l'énoncé et que l'exemple a.}

$A$ explicabilis est de qualibet illarum quatuor $B.\ D.\ G.\ H.$

$50N - 35Q + 16C - 1QQ$ æquatur $24$

fit $1N$ \struck{\ill{}} $1.\ 2.\ 3.$ vel $4$.
\marginal{10}
\note{« 16 » est souligné, et « 10 », souligné, est écrit en regard dans la marge de gauche : c'est la correction, $(x-1)(x-2)(x-3)(x-4) = x^4 - 10x^3 + 35x^2 - 50x + 24$. Devant « 1. », un signe biffé.}

\subsection*{prop. IIII}

\[ \text{Si } \left.\begin{array}{l} A\ \text{quadrato-cubus} \\[1ex] -\left\{\begin{array}{l} B \\ D \\ G \\ H \\ K \end{array}\right\} \text{in } A\ \text{quadrato-quad.} \\[1ex] +\left\{\begin{array}{l} B \text{ in } D \\ B \text{ in } G \\ B \text{ in } H \\ B \text{ in } K \\ D \text{ in } G \\ D \text{ in } H \\ D \text{ in } K \\ G \text{ in } H \\ G \text{ in } K \\ H \text{ in } K \end{array}\right\} \text{in } A\ \text{cubum.} \end{array}\right. \]
\note{Les dix produits sont rangés en trois groupes séparés par des traits (ceux de $B$, ceux de $D$, puis $G$ in $H$, $G$ in $K$, $H$ in $K$). L'accolade de droite descend jusqu'au pied du feuillet sans second membre : l'énoncé continue au-delà de la vue 140.}

\end{document}
